% 1. Preamble
\documentclass[12pt, a4paper]{article}

% Import Packages
\usepackage[utf8]{inputenc}        % Encoding
\usepackage{amsmath, amssymb}     % Advanced math formulas
\usepackage{mathtools}           % Math enhancements
\usepackage{graphicx}            % Image insertion
\usepackage{hyperref}            % Interactive hyperlinks & PDF bookmarks
\usepackage{color}               % Color text and backgrounds
\usepackage{geometry}            % Page geometry
\usepackage{comment}             % Multi-line comments
\usepackage{titling}             % Custom title formatting

\geometry{a4paper, margin=1in}

% Custom Command Definitions
\newcommand{\Sball}{\mathcal{S}}
\providecommand{\norm}[1]{\left\lVert#1\right\rVert}

% Hyperref Setup for Clean PDF Links
\hypersetup{
    colorlinks=true,
    linkcolor=blue,
    urlcolor=blue,
    pdftitle={Pattern Arithmetic},
    pdfauthor={Sandeep Lakshminarayan Chiluveru}
}

% Metadata & Custom Title Banner
\pretitle{\begin{center}\Huge\fontfamily{tt}\selectfont ================================================================================\\[0.5em]}
\posttitle{\\[0.5em] ================================================================================\end{center}}

\title{Pattern Arithmetic: Typed Algebra of Bounded Units on a Phase-Colored Sphere}
%\title{Pattern Arithmetic: Bounded High-Density Units on a Phase-Colored Sphere}
%\title{Pattern Arithmetic: Typed Algebra of Units on a Phase-Colored Sphere}
%\title{Pattern Arithmetic: Phase-Colored Geometric Framework for High-Density Information via Unit Truth Spheres}
\author{Sandeep Lakshminarayan Chiluveru}
\date{\today}

\begin{document}

\maketitle

% ABSTRACT
\begin{abstract}
Ordinary arithmetic treats the integer 1 as a unit of truth: $1=1$ and $1+1=2$. Counting is only one algebra. Colors and musical notes are also units of truth; combining them yields a hue or a chord, not a larger integer. This paper treats a pattern as a higher-order unit and represents one pattern by one typed point
$x = (r,\,\hat n,\,\chi,\,\tau)$
on a unit sphere $\mathcal{S} \subset \mathbb{R}^3$\footnote{Three
dimensions is the standing assumption throughout --- every axiom,
proof, and worked example below uses $\mathbb{R}^3$ unless a section
states otherwise. Section~\ref{subsec:transform_higher_dim} extends
$T$ to $S^{N_{\mathrm{dim}}-1}$ for $N_{\mathrm{dim}}\ne3$; that is a
named exception, not a change to the default, as is the
general-dimension form of offset's transport in
\S\ref{subsec:offset_logexp}.}:
strength $r \in [0, 1]$,
direction $\hat n \in S^2$ (equivalently $(\theta,\phi)$),
internal phase-color $\chi \in [0, 2\pi)$,
and sort $\tau$.
Pattern Arithmetic is the algebra of such units. Five verbs answer
five questions. Gather ($\oplus$) lists who is present and does not
double identical states ($A \oplus A = A$); a gathered unit may be
removed ($\{A\} \ominus A = \emptyset$). Mix ($\boxplus$) reads a
composite; anti-phase cancellation to black is that reading, not a
gather. Transform ($T$) walks an order so the same ingredients may
form different expressions. Offset reuses a walk. Assembly
($\bowtie$) relates two units without fusing them --- the first
operator here that needs anything outside the 4-tuple: an external
seat. Units left alone descend an energy until a stated threshold
freezes them, resting in one camp or several; left longer, all would
merge into one. Collapse~$\to$~$T$~$\to$~result is
one pipeline; assemble~$\to$~rest is the other. Division as
unmixing is left open. Words and numbers occupy the same sphere under different encodings and meet in a document by order and typed action, not by blending. The aim is a single bounded register that carries overlapping identities, structural paths, and mixed prose without a discrete grid. The sphere and its operators are one proposed encoding of a narrower claim --- a combination rule should match the unit --- not a theory of language, physics, biology, or computation.
\end{abstract}

\section{Introduction}

Ordinary arithmetic uses the integer $1$ as a unit of truth.
The identity $1=1$ says that two tokens are of the same kind.
The rule $1+1=2$ says that combining two such tokens produces a
larger quantitative state. That algebra answers the question
\emph{how many?}

It is not the only algebra. A color and a musical note are also
units of truth, but their combination does not enlarge a pile.
Under additive light, red and green yield yellow. Two copies of
the same yellow, once brightness is normalized, remain yellow.
Idempotence is a property of $\oplus$ alone: $A\oplus A=A$ always,
whereas $A\boxplus A$ grows strength to $r'=\min(1,2r)$ at unchanged
hue, and returns the same yellow only after that normalization
(\hyperref[subsec:axiom_interference]{Axiom~5}).
Two notes yield a chord, not a larger integer; a tone plus its
opposite phase yields silence. In each case the word ``add''
names a different rule, and therefore a different kind of unit.

\begin{quote}
\textbf{Thesis (an arithmetic of units of truth).}
\emph{Mathematics need not be performed on numbers alone. Wherever a
domain supplies a unit of truth --- the smallest thing it treats as
one --- and a small set of rules by which such units combine, obeyed
consistently within that domain, the unit and its rules together form
an arithmetic of that domain. Counting is one instance: the unit is a
token and the rule is $+$. Such rules are seldom invented. They are
discovered, as the outcomes a domain returns every time the same
conditions are repeated and it is left to settle toward stability.}
\end{quote}

Examples are everywhere once looked for. Light combines additively
and pigment subtractively: one domain, color, carries two
arithmetics, because it is the rule set and not the domain that
defines one. Notes combine into chords; a plate driven at the same
frequency settles into the same nodal figure
(\S\ref{subsec:chladni_inclusion}); water freezing under the same
conditions returns the same six-fold symmetry. The oldest instances
are older than any notation for them. Atoms combine only in the
proportions their valences allow, with mass and charge conserved
across the reaction. Conserved quantities in particle physics,
electric charge among them, add across every interaction, and some
combinations are forbidden outright. In biology, an engulfed
bacterium and its host became one cell, and the mitochondrion still
carries a genome of its own: two units joined into a stable whole
without either being dissolved --- the shape of assembly
(\hyperref[subsec:axiom_assembly]{Axiom~6}), not of a blend.

These domains are cited as evidence that such arithmetics exist, not
as domains this paper models;
\hyperref[subsec:axiom0_admissibility]{Axiom~0} decides which sorts
earn a geometric reading here, and most physical domains do not.
Ice is the plainest case: its six-fold growth is a genuine arithmetic
of its own domain, recorded by Axiom~0 as the snowflake's native law,
and it is exactly the kind refused a map onto this sphere. This
paper states its own rules as axioms rather than discovering them.
The second half of the thesis appears in it only in miniature: in the
enclosure of Section~\ref{sec:enclosure}, units left under one fixed
law settle toward equilibrium and are read at a stated checkpoint.

When the object of interest is a structured whole rather than a
count, a hue, or a pitch, the natural unit is a \emph{pattern}.
This paper represents one pattern by one typed point on a unit
sphere $\mathcal{S} \subset \mathbb{R}^3$, not by a field painted over the globe:
\begin{itemize}
    \item $r \in [0, 1]$ is the strength of the unit, capped so that a state remains a single unit of truth --- an activation, not a global probability density;
    \item $\hat n \in S^2$, or equivalently $(\theta, \phi)$, is the direction of the pattern on the sphere;
    \item $\chi \in [0, 2\pi)$ is an internal phase-color, so two identities may occupy the same direction without overwriting one another;
    \item $\tau$ is the sort ($\mathtt{word}$, $\mathtt{num}$, $\mathtt{color}$, \ldots): a legality tag, not a fifth geometric angle.
\end{itemize}
Pattern Arithmetic is the algebra of these units. Gather does not double identical states ($A \oplus A = A$). A gathered unit may be removed ($\{A\} \ominus A = \emptyset$). Rotations of the sphere preserve the unit. Opposite-phase Mix is a later reading; it is not $\oplus$.

A single numerical picture is enough to see the difference from counting. Take two unit states at the same place, one with phase $0^\circ$ (red) and one with phase $120^\circ$ (green). Counting would return $2$. Phase-aware combination returns a single unit whose blended phase is $60^\circ$ (yellow). That is this paper's color inclusion --- phasor Mix on a toy chart, not CIE or linear RGB (\hyperref[subsec:axiom0_admissibility]{Axiom~0}, \hyperref[subsec:axiom_interference]{Axiom~5}). The same two strengths with opposite phase ($0^\circ$ and $180^\circ$) return the null state, black. Two units of truth may become a new identity, a stronger mark of the same identity, or none.

The working order of the paper is not that mix alone. One pipeline is
\[
\text{Collapse (bowl)} \;\longrightarrow\; T\text{ (path)} \;\longrightarrow\; \text{Result}.
\]
The other is assemble ($\bowtie$) then rest of seats. Both use the
same unit. Mixing is a later reading of a bowl, not a third spine.
Collecting red and green is not yet yellow, and not yet a poem.
Yellow is what you get if the later reading is Mix ($\boxplus$).
Red and green may equally well walk a $T_{\mathrm{poem}}$ path; the
result need not be yellow --- though a given poem may still arrive
there. The later act, not the shells, chooses the result. A verse is
another path.
Each word of ``Twinkle, twinkle, little star\ldots'' is a unit state
on the same sphere. Collapse yields the bag of distinct words
($\texttt{twinkle}$ appears twice in the line and only once in the
bowl). The poem is the path that walks those states in order,
including the repeated $\texttt{twinkle}$. The bag is not the song.
Numericity is not stored by $\oplus$. Two identical units collapse
to one address; that is the point of $A\oplus A=A$. How many times
the address is \emph{visited} lives in the path (the tape says
\texttt{twinkle} twice). How many as a \emph{count} lives in
$\mathtt{num}$ mode, where $1+1+1+1=4$ and four tokens of $1$ are
not collapsed. Strength $r$ is a poor substitute for either: a
second \texttt{twinkle} cannot honestly raise $r$ past $1$, and
``four ones'' is schoolbook addition, not a fuller unit of truth.
The same bowl may be read more than one way. Loading is not meaning.

The sphere is the stage. It is not the deepest claim. Strip away
geometry, phase, rotations, and the language examples, and one
proposition remains: different kinds of information require different
notions of combination, and an operation should be read by what
informational property it \emph{preserves} --- and what it is allowed
to throw away.
\[
1+1=2,
\qquad
\mathrm{red}\boxplus\mathrm{green}=\mathrm{yellow},
\qquad
A\oplus A=A,
\qquad
T_a(T_b(x))\neq T_b(T_a(x)).
\]
These are not alternative tricks for the same plus sign. They answer
different questions:

\begin{center}
\begin{tabular}{ll}
Operation & Question \\
\hline
$+$ & How many? \\
$\oplus$ / $\ominus$ & What is present? / what is removed? \\
$\boxplus$ ($\mathrm{Mix}$) & What composite state results? \\
$T$ & What happens, and in what order? \\
offset & What walk is reused? \\
$\bowtie$ & How do two units relate while both still sit? \\
\end{tabular}
\end{center}

Division as deconvolution or signal unmixing is not among these
verbs; it is left open
(Section~\ref{subsec:division_open}).
Inverse rotation of $T$ is well-defined. A Mix inverse is not.

A pattern can be treated as one unit. Identity and multiplicity come
apart: $\oplus$ keeps kind and discards count; the path keeps visits;
$\mathtt{num}$ keeps quantity. Composition and transformation are
different verbs. Sorts may share a stage without sharing an operator.
The geometry that follows is only a language in which those verbs can
be written. It does not mint the relations; a lexicon or a learning
rule must still seat $\mathrm{Sun}$ and $\mathrm{Day}$.
Section~\ref{subsec:claim_scope} keeps that claim narrow.

\subsection{The Limitations of Discretized Matrix Grids}
A discrete matrix grid is the wrong box for such a unit. Traditional matrix representations of continuous state spaces suffer from two severe structural bottlenecks:
\begin{itemize}
    \item \textbf{Memory Scaling Bottlenecks:} A voxel cube of resolution $N_{\mathrm{vox}}$ stores $O(N_{\mathrm{vox}}^3)$ cells, growing steeply with resolution.
    \item \textbf{Signal Identity Destruction:} Mapping multiple overlapping variables to identical spatial coordinates often results in the destructive overwriting of underlying signal identities.
\end{itemize}

\subsection{The Phase-Colored Unit Sphere Framework}
To resolve these grid constraints, the sphere is offered as one bounded object that can carry overlapping identities by direction and phase, without enlarging a lattice.
\begin{itemize}
    \item \textbf{Spatial Orientation $\hat n \in S^2$, or $(\theta, \phi)$:} Defines continuous directional coordinates across the spherical surface.
    \item \textbf{Magnitude ($0 \leq r \leq 1$):} Strength of that one unit, capped so the arrow remains a single unit of truth. An activation, not a density for the whole sphere.
    \item \textbf{Internal Color Phase ($\chi \in [0, 2\pi)$):} Serves as an independent phase degree of freedom.
    \item \textbf{Sort $\tau$:} A legality tag ($\mathtt{word}$, $\mathtt{num}$, $\mathtt{color}$, \ldots). It does not move the arrow.
\end{itemize}

\subsection{Core Contributions \& High-Level Intuition}
The principal contributions of Pattern Arithmetic are:
\begin{itemize}
    \item \textbf{A typed unit:} one atom $x=(r,\hat n,\chi,\tau)$ on a bounded sphere, not a voxel cell and not a global probability field.
    \item \textbf{Admissibility of sorts (Axiom~0):} a sort earns a geometric $\chi$- or $T$-reading only where a stated map preserves its native combination; most domains do not.
    \item \textbf{Verbs on that unit:} $\oplus$ collects identities ($A\oplus A=A$); $\boxplus$ reads a composite; $T$ walks an ordered path on one register occupant; offset copies a relation; $\bowtie$ assembles two units into a relation without fusing them. Schoolbook $+$ is reserved for $\mathtt{num}$.
\end{itemize}

Later sections distinguish a numeric encoding from a lexical encoding on the same sphere, so that a report may carry words and numbers together without forcing them through $\oplus$.

\subsection{Scope of the claim}
\label{subsec:claim_scope}
The actual claim is narrow. Counting ($1+1=2$) is not the only valid
arithmetic for combining units of information. A combination rule
should be chosen for what the unit \emph{is} --- a count, a color, a
direction, a word --- not applied by default because it is the
arithmetic everyone already knows. The sphere, phase, and the
operators $\oplus$, $\boxplus$, $T$, offset, and $\bowtie$ are one
proposed way to make that claim precise and checkable. They are not
themselves the claim, and they are not offered as a general theory of
language, physics, biology, or computation.
The paper does not attempt to compute the governing equations of
those systems. It attempts to arrive at higher-order realities by
reading many co-present units of truth --- seated units, whether
pinned by a lexicon or by a learned map --- through a path, or by
letting external seats rest under assembly. That is a computation of
\emph{presence, phase, order, and rest}, not a substitute for a field
equation. The enclosure's energy is a bookkeeping rule for seats
$\vec P$; it is not a model of gravity or of matter.
Storms, forests, a table and chair, and a pathogen cancellation later in the paper
(Section~\ref{sec:higher_order}) are costumes for the
$\oplus$-versus-$T$ split: presence versus path. Related-work
paragraphs (Section~\ref{sec:related}) draw boundaries against
existing formalisms. None of that is a demonstrated model of weather,
ecology, or epidemiology, or a replacement for quantum mechanics,
embeddings, or transformers. A later widening of a costume into a
domain claim would need its own evidence. It is not an extension of
the illustration.
Section~\ref{subsec:axiom0_admissibility} states this limit formally:
a sort earns a geometric reading only where a structure-preserving
map into Mix or $T$ exists at a stated fidelity, and most physical
domains do not offer one.

Every law in this paper is deterministic: the same starting units
give the same result, step for step. Where a claim rests on a
starting configuration, the configuration itself is printed, as in
the worked examples of Sections~\ref{sec:transformation_algebra}
and~\ref{sec:higher_order} and the enclosure of
\S\ref{subsec:enclosure_replay}. How a reader generates
configurations of their own --- by hand, or by any library under any
seed --- is left open; the printed one is the shared base. Printing
it makes a result replayable. It does not make the result right.

The next section turns these definitions into formal mathematical axioms.


\section{Axiomatic Foundations Of Pattern Arithmetic}

\subsection{Axiom 0: Admissibility of Sorts}
\label{subsec:axiom0_admissibility}
The axioms that follow govern how a state $x=(r,\hat n,\chi,\tau)$
behaves once it sits on $\mathcal{S}$. They say nothing about which
sorts $\tau$ have earned the right to sit there with a geometric
reading attached. That prior question is this axiom, and it is
separate from Table~\ref{tab:type_table} in
Section~\ref{subsec:type_table}: the table says which
\emph{operators} two already-admitted sorts may use on each other;
this axiom says which sorts may claim $\chi$ and $T$ mean something
about their own domain in the first place.

\paragraph{The test.} A sort $\tau$ has a domain $D_\tau$ and, in the
world, its own native way of combining two members of that domain,
written $a \circ_\tau b$ for $a,b \in D_\tau$. Call $\tau$
\emph{geometrically admissible} if there is a map
\begin{equation}
    \varphi_\tau : D_\tau \longrightarrow \mathcal{S}, \qquad
    \varphi_\tau(a \circ_\tau b) \;\approx\;
    \boxplus\bigl(\varphi_\tau(a), \varphi_\tau(b)\bigr)
    \quad \text{or} \quad
    T_{\varphi_\tau(a)}\bigl(\varphi_\tau(b)\bigr),
    \label{eq:admissibility}
\end{equation}
holding at a stated fidelity (exact, or approximate with the
approximation named). Equation~\eqref{eq:admissibility} is the
Tier~1 test. Tier~0 admission, below, does not require it.
If no such $\varphi_\tau$ exists, $\tau$ is
\emph{geometrically inadmissible} for Tier~1: it may still occupy a point on
$\mathcal{S}$, but only under the inert treatment described below,
never with a claim that $\chi$ or $T$ tracks anything the domain
itself does.

\paragraph{Three tiers.}
\begin{itemize}
    \item \textbf{Tier 0 --- default, unstated.} $\mathtt{word}$ and
    $\mathtt{num}$ need no admissibility argument because neither
    makes a physical claim through the geometry. $\mathtt{num}$
    has an exact, invertible \emph{chart}
    (Section~\ref{subsec:clock_encoding}): decode, add in
    $\mathbb{Z}$, encode. That chart is not a Mix or $T$
    homomorphism --- Mix on it is the clock trap of that same
    section, which is why Mix on $(\mathtt{num},\mathtt{num})$ is
    type-refused. $\mathtt{word}$ makes no claim at all --- $T$ on a
    word is bookkeeping (order, presence), never a statement about
    linguistic physics. Mix applied to two word-states is inert
    geometry, not a $\chi$-reading of language. Both are silently
    admitted at this inert level; nothing further needs to be said
    to use them.

    \item \textbf{Tier 1 --- explicit inclusion.} A sort earns a
    genuine $\chi$-reading only by stating its $\varphi_\tau$ and the
    fidelity it holds at. Musical pitch class is close to exact:
    octave equivalence makes pitch native to a circle, so $\chi$ for
    a note is close to literal, not borrowed. Color is close but not
    unique: additive light mixing is linear in some models and only
    approximately phasor-like in others, so an inclusion of
    $\mathtt{color}$ must say \emph{which} reading it means --- Mix
    under $\boxplus$ and Mix under linear RGB blend are different,
    non-interchangeable inclusions, and a paper that admits color
    must pick one and say so. This paper picks phasor Mix on a toy
    $0^\circ/120^\circ/60^\circ$ chart
    (\hyperref[subsec:axiom_interference]{Axiom~5}). It does not pick CIE,
    and it does not pick linear RGB. A second Tier-1 sort,
    $\mathtt{chladni}$, is worked in full in
    Section~\ref{subsec:chladni_inclusion}, where
    Equation~\eqref{eq:admissibility} is met exactly within a mode
    and shown unavailable across modes.

    \item \textbf{Inadmissible --- no $\varphi_\tau$ exists.} Some
    domains have a native combination law with no structure-preserving
    map to $(\oplus,\boxplus,T,\text{offset},\bowtie)$ at any useful fidelity.
    A snowflake's native symmetry is the discrete point group
    $D_{6h}$, and two ice structures meeting is a diffusion-limited
    accretion process governed by temperature and humidity gradients
    --- not a continuous rotation in $\mathrm{SO}(3)$. There is no
    honest $\varphi_{\mathrm{snowflake}}$ satisfying
    Equation~\eqref{eq:admissibility}. Such a sort is not barred from
    the register outright: it may still be loaded as a point, but
    only demoted to the Tier-0 treatment already given to
    $\mathtt{word}$ --- $T$ used for order or presence only, no
    $\chi$-reading granted, and no claim that a rotation of its point
    corresponds to anything the domain does. What is barred is Tier-1
    status: an inadmissible sort may never be handed a stated
    $\varphi_\tau$, because none would be true.
\end{itemize}

\begin{table}[h]
\centering
\begin{tabular}{llcc}
\hline
Sort & Native law $\circ_\tau$ & Mix/$T$ map? & Tier \\
\hline
$\mathtt{num}$   & schoolbook $+$ & no (chart only) & 0 \\
$\mathtt{word}$  & none claimed & n/a & 0 \\
$\mathtt{note\text{-}pitch}$ & pitch-class sum & near-exact & 1 \\
$\mathtt{color}$ & additive light & phasor Mix (here) & 1 \\
$\mathtt{chladni}$ & modal superposition & exact, within a mode & 1$^{\S}$ \\
$\mathtt{snowflake}$ & $D_{6h}$ growth & no & T1 no; T0 yes \\
\hline
\end{tabular}
\caption{Admissibility is a property of the sort, decided before
Table~\ref{tab:type_table}'s operator legality is even asked.
Equation~\eqref{eq:admissibility} is the Tier~1 test; Tier~0
admission does not use it. A sort
absent from this table is undecided, not admitted; treat it as
inadmissible until a $\varphi_\tau$ is stated.
$^{\S}$$\mathtt{chladni}$ is admitted at Tier~1 with a stated
restriction: exact for superposition within one mode, undefined across
modes (Section~\ref{subsec:chladni_inclusion}).}
\label{tab:admissibility}
\end{table}

This axiom is also what gives
Section~\ref{subsec:claim_scope} teeth. ``Not a general theory of
physics, biology, or computation'' is not only a promise of
restraint; it is a consequence of Equation~\eqref{eq:admissibility}
failing for most physical domains a reader might otherwise expect
the sphere to cover.

\subsection{Process principle: Collapse, path, result}
\label{subsec:process_principle}
The $A_i$ below are already seated units, each loaded from an
observed sequence by $E_L$ (Equation~\eqref{eq:lexicon_cast});
collapse is the act that discards the order and count of that
sequence, not the act that produces the units themselves. With that
understood, Pattern Arithmetic proceeds in three named acts, kept separate so that a bowl may grow before it is read:
\begin{equation}
    \underbrace{A_1 \oplus A_2 \oplus \cdots \oplus A_n}_{R\ =\ \text{the bowl}}
    \;\longrightarrow\;
    T_{w_k}\circ \cdots \circ T_{w_1}(x)
    \;\longrightarrow\;
    \text{Result}.
    \label{eq:pipeline}
\end{equation}
Here $x$ is the register occupant; the $w_i$ are actors, taken from
the bowl or from the tape. $T$ does not act on the set $R$ as a set.
$\oplus$ \emph{collects}: identical tokens do not double ($A \oplus A = A$), and opposite marks may sit together until a later reading. $T$ is an ordered path, not a mix. The result is whatever that path (or another named reading) makes of the bowl --- a poem, a report, a staged endpoint.

Paint is Mix ($\boxplus$), a reading of peer states at a common direction, not a $T$:
\[
\mathrm{red}\boxplus\mathrm{green} = \mathrm{yellow}.
\]
The same pair may take $T_{\mathrm{poem}}$ instead. That result is a
verse about the two colors, not a hue --- unless the poem itself
walks back to yellow. Shells do not own a pathway. A verse is another
path. One bowl, several ways to read it. One may collect
and read in a single gesture; the algebra does not force the merge.
This principle is an axiom of \emph{order of acts}, not a new formula
for a 4-tuple.

\paragraph{Levels are closed.}
A result may be seated again. Equation~\eqref{eq:pipeline} returns a
state, and a state may take a seat in a further bowl, so the three
acts repeat at a higher level: words gather and are walked into a
line; the line is seated; lines gather and are walked into a verse.
Two properties hold at every level and are worth stating once. First,
a unit carries no representation of what it constitutes. A seated word
does not know the poem, and $T$ leaves $\tau_x$ invariant
(Section~\ref{subsec:transform_properties}), so being acted upon does
not change what a unit is. Second, a bowl carries no representation of
the reading that will be applied to it; that is what \emph{one bowl,
several ways to read it} means above.

The consequence is a design rule rather than a theorem. Whatever a
lower unit would have to know about a higher one belongs at the higher
level and not in the 4-tuple: which units were taken as co-present,
over what interval, within what region, and under what question. Those
are constituting facts of a level, not coordinates of an atom, and
keeping them out is why the atom stays at four coordinates and one
sort tag while the levels above it may be as elaborate as a problem
requires. This paper works a single level and states the mechanism for
the next; it does not carry a multi-level example.

\subsection{Axiom 1: The Normalized Unit Truth State ($r \le 1$)}
\label{subsec:axiom1}
A pattern is one typed point, not a field painted over the globe. The working atom is
\begin{equation}
    x = (r,\,\hat n,\,\chi,\,\tau),
    \qquad
    \begin{aligned}
    r &\in [0,1],
    \qquad
    \hat n \in S^2,\\
    \chi &\in [0,2\pi),
    \qquad
    \tau \in \{\mathtt{word},\,\mathtt{num},\,\mathtt{color},\,\mathtt{note\text{-}pitch},\,\mathtt{chladni},\,\ldots\}.
    \end{aligned}
    \label{eq:unit_truth}
\end{equation}
Equivalently, $\hat n$ may be written in polar angles $(\theta,\phi)$, so the geometric part of the atom is the 4-tuple $(r,\theta,\phi,\chi)$ used by the operators below. The sort $\tau$ does not move the arrow; it says which buttons may be pressed (Section~\ref{sec:encodings_sorts}). A collection (bowl) is a separate finite set $B=\{x_1,\ldots,x_n\}$ of such atoms. The sphere holds one current occupant of the register; the bowl is who is present.

The teaching atom is this 4-tuple: one arrow, $N=1$. \hyperref[subsec:axiom_v1_metrics]{Axiom~7} allows an atom to hold $N$ such arrows under one closure, each still capped at $r_i\le 1$, plus one uncapped derived sum that is not itself a unit. That bundle does not replace Equation~\eqref{eq:unit_truth}. Every informational arrow in the bundle is an instance of this axiom. Operators that need a single pair $(r,\hat n)$ read the derived sum as a weight or a direction; as the actor of $T$ in three dimensions, only after the cast of \hyperref[subsec:axiom_v1_metrics]{Axiom~7}, and beyond three dimensions not from the derived sum at all, but from $N_{\mathrm{dim}}-2$ of the bundle's own orthonormal arrows directly (\S\ref{subsec:transform_higher_dim}). They do not mint a second kind of unit.

The cap $r \le 1$ applies to \emph{that} arrow. It is a normalized strength / activation coordinate, not a global probability density and not a demand that all arrows on $\mathcal{S}$ integrate to one. Two states may occupy
different directions at once; each still carries its own strength
$r\le 1$. The sphere is the stage. The unit is the vector.
$r=1$ is a full-strength mark. Nothing in this axiom sums the globe
into a single global budget. A field of many atoms, if wanted later, is a different object (for example a superposition of Dirac masses at those points), not the atom itself.

\paragraph{How the bound is imposed is a choice.}
This paper bounds by clipping: Mix returns $r'=\min(1,|z'|)$, so
$r=1$ is attained whenever a phasor sum reaches unit length. Clipping
is not injective. Every sum of length one or more returns the same
$r'$, and agreement past that point is not recorded
(Section~\ref{subsec:axiom_interference}). An order-preserving
bijection from $[0,\infty)$ onto $[0,1)$ --- $|z|\mapsto|z|/(1+|z|)$,
say --- would bound without discarding, since distinct sums would keep
distinct strengths, and it would leave $z=0\mapsto 0$ and the black of
\hyperref[subsec:axiom_null]{Axiom~2} untouched.

The clip is kept here, and the trade is stated rather than argued.
Under such a bijection $r=1$ becomes a limit and not a value: the
saturated pole of \hyperref[subsec:axiom_white]{Axiom~3} would be approached
and never occupied, and $\alpha=\pi r_w$ would never reach the
half-turn of Section~\ref{subsec:transform_properties}. Two attained
poles are worth more to this paper than lossless accumulation, because
no construction here sums a large bowl. A setting that does --- a
verdict accumulated from many units, say --- would weigh the same
trade the other way, and should revisit this line rather than inherit
it.

\subsection{Axiom 2: The Absolute Null State ($r = 0$, Black)}
\label{subsec:axiom_null}
The lower bound of a single unit is vanishing strength:
\begin{equation}
    r = 0 \quad \implies \quad x = \mathbf{0}
    \quad \text{(Black)}.
    \label{eq:null_state}
\end{equation}
Direction and phase are then unused: there is no arrow to point.
Black is an empty seat, not a field that is zero at every angle.

\subsection{Axiom 3: Saturation (White)}
\label{subsec:axiom_white}
White is the upper pole of the same bounded stage. It is not an infinite
norm. Two readings are allowed; neither leaves the cap $r\le 1$.

\paragraph{White as a colour.}
A single arrow may take a full-spectrum name (white light as one
identity). That state is still one 4-tuple with $r\le 1$.

\paragraph{White as a full house.}
The bowl may be saturated: many directions occupied, many hues present.
Collecting further distinct shells fills more seats; it does not raise
any arrow past $r=1$. $\oplus$ does not grow strength. Mix may grow a
single $r$ until the cap. The stage looks white because the palette is on. The catalogue
of names is large --- on the order of a 24-bit palette,
$\sim 16$ million colours per direction --- and may be called
practically endless. No unit is $r=\infty$.
\begin{equation}
    \text{White} \;=\; \text{saturation of the bowl, or a full-spectrum name.}
    \label{eq:white_state}
\end{equation}
Black is empty. White is as full as this construction allows. The
menu is huge. The arrows stay finite.
Density in this paper means many identities can share the stage
without overwrite --- many directions, and many phases per direction
--- each still capped at $r\le 1$. It does not mean infinite
information, and it is not a storage theorem against a voxel grid.

One atom carries a modest, bounded amount before it runs out of room.
Two routes exist for more, and they are not the same route. The first
is to seat a second atom: $\oplus$ gathers it into the same bowl
(Section~\ref{subsec:axiom_idempotence}), which is how this
construction represents anything beyond what one vector alone can
hold, and is the route already in use throughout this paper. The
second is to raise the ambient dimension, $S^2\to S^{N_{\mathrm{dim}}-1}$
(Section~\ref{subsec:dimensionality_open}), giving a single atom more room
on the sphere it already occupies. The second route is conditional on
the first: it relieves crowding among many co-present atoms, and a
bowl too small to crowd gains nothing by it --- two atoms do not
need a millionfold larger stage to stay apart.

\subsection{Axiom 4: Gather and remove ($\oplus$, $\ominus$)}
\label{subsec:axiom_idempotence}
$\oplus$ \emph{collects} a unit into the bowl. It does not mix colours
and it does not turn an arrow. A bowl is a finite subset of
$\mathcal{S}$, and $\oplus$ is union on such subsets:
\[
\oplus:\; 2^{\mathcal{S}}\times 2^{\mathcal{S}}
\longrightarrow 2^{\mathcal{S}},
\qquad
B_1 \oplus B_2 \;=\; B_1 \cup B_2 .
\]
A bare state $A$ appearing in an $\oplus$-expression denotes its
singleton bowl $\{A\}$. With that convention, identical names occupy
one seat:
\begin{equation}
    A \oplus A = A,
    \qquad\text{that is,}\qquad
    \{A\}\cup\{A\}=\{A\},
    \label{eq:idempotence_formal}
\end{equation}
while distinct names take separate seats, $A\oplus B=\{A,B\}$, so
opposite marks may sit together:
$\{+1\}\oplus\{-1\}=\{+1,-1\}$.
Idempotence, commutativity and associativity of $\oplus$ are those of
union, and none of the three needs separate postulation.
Gather is sort-blind: any two seated units may share a bowl, whatever
their sorts, since a bowl records presence and presence makes no claim
about a domain (Table~\ref{tab:type_table}). Sort discipline governs
what may later be \emph{done} with the bowl, not who may sit in it.
The bowl is a set of states, not a blended 4-tuple.

$\ominus$ \emph{removes} a unit,
$\ominus:\,2^{\mathcal{S}}\times\mathcal{S}\to 2^{\mathcal{S}}$.
Removing the last occupant leaves the empty bowl:
\begin{equation}
    \{A\}\ominus A = \emptyset.
    \label{eq:ominus_empty}
\end{equation}
The empty bowl $\emptyset$ and the null state $\mathbf{0}$ of
\hyperref[subsec:axiom_null]{Axiom~2} are distinct objects: $\emptyset$ is a
bowl with no occupant, $\mathbf{0}$ is an occupant with $r=0$. Only
$\mathbf{0}$ is a state, which is why $T_{\mathbf{0}}$ is defined
(Section~\ref{subsec:transform_properties}) and $T_{\emptyset}$ is
not.
$\ominus$ undoes a gather. It is not anti-phase interference.
Idempotence then gives the counting-trap
\begin{equation}
    A \oplus A \oplus A \oplus A \ominus A = \emptyset
    \quad\text{(the empty bowl)}.
    \label{eq:four_a_ominus}
\end{equation}
Four gathers of the same name are still one seat; one removal empties
it. Schoolbook $1+1+1+1-1=3$ lives only in $\mathtt{num}$ mode.

\subsection{Axiom 5: Mix as a reading (\texorpdfstring{$\boxplus$}{boxplus})}
\label{subsec:axiom_interference}
A blended 4-tuple appears only when Mix \emph{reads} the bowl.
Write $\mathrm{Mix}=\boxplus$. Mix is a \emph{phasor} operator: it
combines $\chi$ among states that already share a direction $\hat n$.
Direction belongs to $T$ and to the rotations of
\S\ref{subsec:axiom_rotation}; Mix never moves an arrow onto a new
axis. Mix of states with differing $\hat n$ is left open
(Section~\ref{subsec:mix_scattered_open}). The output is
$(r',\theta',\phi',\chi')$ with $r'\le 1$, or black.
$T_w$ is a one-argument rotation and phase-shift of an occupant,
parameterized by an actor $w$. Mix is a combination of
peer states. It is not a member of the group $T_w$ lives in, and it
does not compose with $T_w$ the way two ordinary $T$'s do. The two
give different answers on the same pair: for unit red at $\chi=0^\circ$
and unit green at $\chi=120^\circ$ sharing a direction,
$T_{\mathrm{green}}(\mathrm{red})$ shifts phase additively to
$\chi'=120^\circ$, while $\mathrm{red}\boxplus\mathrm{green}$ takes the
phasor resultant, $\chi'=60^\circ$. No choice of actor makes $T$
return $60^\circ$ here, because $T$ adds phase where Mix sums phasors.
Three further properties separate them: $T_w$ is asymmetric in actor
and patient while $\boxplus$ is commutative; $T_w$ leaves $r$ invariant
while $\boxplus$ may grow it to the cap
(\hyperref[subsec:axiom_white]{Axiom~3}); and every $T_w$ has an inverse
$w^{*}$ while $\boxplus$ has none.

Which sorts may read a Mix result as a claim about their own domain is
settled by \hyperref[subsec:axiom0_admissibility]{Axiom~0}, not here: Tier~1
for $\mathtt{color}$ and $\mathtt{note\text{-}pitch}$, inert geometry
for $\mathtt{word}$, type-refused for $\mathtt{num}$.

\paragraph{Two states at a common direction.}
Treat strength and phase as a phasor $z=r e^{i\chi}$. A domain reading
requires $A$ and $B$ to share a sort (\S\ref{subsec:axiom0_admissibility}
above), so the sort carries through unchanged:
\begin{equation}
    z' = r_A e^{i\chi_A} + r_B e^{i\chi_B},
    \qquad
    r' = \min(1,|z'|),
    \qquad
    \chi' = \arg(z'),
    \qquad
    \tau' = \tau_A = \tau_B,
    \label{eq:interference_formal}
\end{equation}
and $(\theta',\phi')=(\theta_A,\phi_A)$. If $z'=0$, the result is
black. Equal unit paints at $0^\circ$ and $120^\circ$ give
$z'=e^{i\pi/3}$, so
\[
(1\angle 0^\circ)\boxplus(1\angle 120^\circ)
=(1,\,\theta,\,\phi,\,60^\circ,\,\mathtt{color})
\quad\text{(yellow)}.
\]
This paper's color inclusion
(\hyperref[subsec:axiom0_admissibility]{Axiom~0}, Tier~1) is that phasor Mix
on a toy $0^\circ/120^\circ/60^\circ$ chart. It is not CIE, and it is
not linear RGB.
The unequal pair $r=0.6$ at $0^\circ$ and $r=0.8$ at $120^\circ$
gives $z'=0.20+0.6928\,i$, hence
\[
r'\approx 0.721,\qquad \chi'\approx 73.9^\circ.
\]
Equal anti-phase marks, same place, give $z'=0$ (black).
A state mixed with itself at the same phase is not idempotent:
$A\boxplus A$ gives $z'=2r e^{i\chi}$, so $r'=\min(1,2r)$.
Strength grows until the cap clips it. That is distinct from
$A\oplus A=A$.

\paragraph{Reading a whole bowl.}
Mix over a bowl $B$ whose members share a direction $\hat n$ and, for
a domain reading, a single sort $\tau$ is
defined in one step, not as a fold:
\begin{equation}
    \boxplus(B)
    \;=\;
    \bigl(\min(1,|z_B|),\;\hat n,\;\arg z_B,\;\tau\bigr),
    \qquad
    z_B \;=\; \sum_{x\in B} r_x e^{i\chi_x}.
    \label{eq:mix_bowl}
\end{equation}
The binary $\boxplus$ of Equation~\eqref{eq:interference_formal} is
the case $|B|=2$. The one-step form is necessary, not stylistic,
because the cap makes the binary operator non-associative: with
$A=1\angle 0^\circ$, $B=1\angle 0^\circ$, $C=1\angle 180^\circ$,
\[
(A\boxplus B)\boxplus C=\text{black},
\qquad
A\boxplus(B\boxplus C)=1\angle 0^\circ ,
\]
since $A\boxplus B$ is clipped from $2\angle 0^\circ$ down to
$1\angle 0^\circ$, and the strength discarded by that clip is exactly
what $C$ would otherwise have cancelled.
Equation~\eqref{eq:mix_bowl} on the same three states sums first and
caps once:
\[
z_B = e^{i0}+e^{i0}+e^{i\pi} = 1,
\qquad
\boxplus\bigl(\{A,B,C\}\bigr) = (1,\,\hat n,\,0^\circ,\,\tau),
\]
one reading, reached without choosing an order. $\boxplus$ is
commutative,
and Equation~\eqref{eq:mix_bowl} is order-independent, so a bowl ---
which carries no order --- has exactly one Mix reading. Iterated
binary Mix does not, and is not what
Equation~\eqref{eq:mix_bowl} means. Non-associativity is the price of
the cap of \hyperref[subsec:axiom1]{Axiom~1}: without $\min(1,\cdot)$ the
phasor sum would be ordinary complex addition and associativity would
be free.

$\oplus$ does not run this rule. Two shells in a bowl have no
resultant angle until Mix (or a $T$-path) is applied.
A reading returns a state and leaves the bowl intact: $\boxplus(B)$
does not consume $B$, and a $T$-path draws actors from a bowl without
removing them. Only $\ominus$ takes a member out. That is what makes
one bowl readable more than once --- the same $B$ may yield a Mix
result and, separately, a walked path.
Gather is lossless and is undone by $\ominus$; Mix is lossy, since it
discards both order and any strength above the cap. That is why
$\ominus$ inverts $\oplus$, $w^{*}$ inverts $T_w$, and nothing inverts
$\boxplus$.

\paragraph{Saturation and levels.}
A sum carries no bound of its own; the cap of
\hyperref[subsec:axiom1]{Axiom~1} supplies one. A bowl of many agreeing units
therefore reaches $r'=1$ quickly, and further agreement past that
point is not recorded. This is worth reading as a remark about levels
rather than a defect in the cap. A phasor sum is a reading of units
that stand at one level together, and repeated samples of a single
unit do not stand at one level; a bowl that saturates early is often a
level that has been flattened
(Section~\ref{subsec:process_principle}). Strength does not
accumulate across levels: a reading returns a state whose $r$ is legal
by construction --- $\min(1,|z_B|)$ for Mix, and the occupant's own
$r$ for a path --- and that state is seated afresh at the level above.
Composition hides the question, since a path over many actors stays in
$\mathrm{SO}(3)$ and never saturates, so it is Mix rather than the
walk that makes a flattened level visible. What the cap does not
supply is a way to average, which needs the count that $\oplus$
discards (Section~\ref{subsec:axiom_idempotence}); reducing repeated
samples to one seated unit is therefore work for the level below, not
for $\boxplus$.

\subsection{A worked Tier-1 inclusion: \texorpdfstring{$\mathtt{chladni}$}{chladni}}
\label{subsec:chladni_inclusion}
Color is admitted on a declared toy chart. A second sort is admitted
here on stronger ground, to show that
Equation~\eqref{eq:admissibility} can be met exactly rather than
approximately, and by physics unrelated to either color or language.

\paragraph{Native law.}
A thin plate driven at one frequency responds as a superposition of
its eigenmodes $\varphi_j(x,y)$, each carrying a complex amplitude
$A_j=a_j e^{i\chi_j}$; the visible figure is the nodal set of
$w=\sum_j A_j\varphi_j$. So $\circ_{\mathtt{chladni}}$ is complex
addition of mode amplitudes. Two drives on the same mode add as
complex numbers.

\paragraph{Chart.}
One state is one mode. Let $\hat n$ carry the mode index under any
injective map $(m,n)\mapsto(\theta,\phi)$, let $\chi$ be that mode's
phase, and let $r=a/a_{\max}$ be its amplitude normalized to a
declared saturation $a_{\max}$.

\paragraph{The map, and where it is exact.}
Within one mode, $\hat n$ is fixed, and
\begin{equation}
    \varphi_{\mathtt{chladni}}(A\circ B)
    \;=\;
    \varphi_{\mathtt{chladni}}(A)\;\boxplus\;\varphi_{\mathtt{chladni}}(B)
    \label{eq:chladni_hom}
\end{equation}
holds identically while $|z|\le 1$, because complex addition of
amplitudes and phasor addition of $(r,\chi)$ are the same operation.
Nothing is fitted. Note also that the chart's arbitrariness costs
nothing here: same-mode combination never moves $\hat n$, so
Equation~\eqref{eq:chladni_hom} places no constraint on the
mode-to-direction map. Two consequences are physical and immediate.
Equal amplitudes in anti-phase give $z=0$: the mode is driven to
silence and its figure disappears, which is
Equation~\eqref{eq:interference_formal} returning black. Two drives at
$0^\circ$ and $120^\circ$ give $r=1$ and $\chi'=60^\circ$ --- the same
arithmetic as $\mathrm{red}\boxplus\mathrm{green}=\mathrm{yellow}$,
on unrelated physics.

\paragraph{The restriction, which is physical.}
Across modes the map is not available, and not for want of effort.
Mode space is not closed under addition: a sum of two distinct modes
is not a mode, so it cannot be seated as a single state under any
chart. Cross-mode coexistence is therefore $\oplus$ --- a bowl of
modes --- and the resulting figure is a reading of that bowl which is
not $\boxplus$. This is the open question of
Section~\ref{subsec:mix_scattered_open} met in a concrete domain, and
it suggests the question may be closed negatively for this sort rather
than left waiting.

\paragraph{Fidelity claimed, and three limits.}
Exact in phase and amplitude below saturation, within one mode;
undefined across modes. Three limits are stated rather than argued
away. The cap $r\le 1$ is a normalization to a declared $a_{\max}$,
not a fact of linear plate theory, so
Equation~\eqref{eq:chladni_hom} clips above saturation, exactly as the
color inclusion does. Direction is a label here and nothing more,
so no meaning is claimed for $T$ on a $\mathtt{chladni}$ state:
rotating a mode index is not an operation the domain performs. And the
figures most often reproduced are degenerate-mode combinations, which
fall on the cross-mode side; what is admitted exactly is the simpler
case, not the famous one.

\subsection{Rotation preserves the unit}
\label{subsec:axiom_rotation}
A rotation $\mathbf{R} \in \mathrm{SO}(3)$ turns the direction $\hat n$ and leaves strength, phase, and sort alone:
\begin{equation}
    \mathbf{R}\cdot(r,\,\hat n,\,\chi,\,\tau)
    \;=\;
    (r,\,\mathbf{R}\hat n,\,\chi,\,\tau),
    \qquad
    \norm{\mathbf{R}\hat n} = \norm{\hat n} = 1.
\end{equation}
The unit of truth is preserved because $r$ is not scaled and $\hat n$ stays on $S^2$. This is the geometric engine of $T$, not a second blend.

\subsection{Axiom 6: Assembly (\texorpdfstring{$\bowtie$}{assembly})}
\label{subsec:axiom_assembly}
Gather lists who is present; Mix fuses peers into one state; $T$
walks an ordered path. A fourth question remains: given two units,
how strongly do they belong together, and where do they touch?
Assembly, written $\bowtie$, answers it by adding a relation while
leaving both units intact --- it is neither $\oplus$ (which states
presence and no relation) nor $\boxplus$ (which fuses and loses both
identities):
\begin{equation}
    A \bowtie B \;=\; (A,\;B,\;\kappa_{AB},\;U_{AB}).
    \label{eq:assembly_signature}
\end{equation}
Both operands survive unchanged. What is new is the pair
$(\kappa_{AB}, U_{AB})$: a contact point and an interaction energy.

\paragraph{The interaction energy.}
\begin{equation}
    U_{AB} \;=\; -\frac{r_A r_B}{2}\cdot
    \frac{3+\hat n_A\cdot\hat n_B}
         {\bigl(d_{AB}^2+\varepsilon_{AB}^2\bigr)^{3/2}},
    \qquad
    \varepsilon_{AB} = r_A + r_B.
    \label{eq:assembly_energy}
\end{equation}
Three factors, each answering a different question. The product
$r_Ar_B$ asks how strong the two units are; a weak partner lowers the
whole. The dot product $\hat n_A\cdot\hat n_B$ asks how well their
directions agree, and is bounded: the bracket ranges over $[2,4]$
only, so alignment can shift the result by at most a factor of two.
The denominator asks how far apart they sit, which requires the
enclosure of Section~\ref{sec:enclosure} and is the one term that
reaches outside the bounded stage of \hyperref[subsec:axiom1]{Axiom~1}. The
softening length $\varepsilon_{AB}$ is derived, not chosen: reading
$r$ as a radius, two units cannot approach closer than the sum of
their own extents.

\paragraph{The contact point.}
\begin{equation}
    \kappa_{AB} \;=\; \hat n_A + \hat n_B,
    \label{eq:assembly_contact}
\end{equation}
the raw, unnormalized sum. Leaving it unnormalized is deliberate.
The normalized bisector $\kappa_{AB}/\lvert\kappa_{AB}\rvert$ is
undefined when $\hat n_B = -\hat n_A$, and that degeneracy drove a
long search for a special case that turned out not to be needed:
$U_{AB}$ never required an axis at all, and $\kappa_{AB}$ is a plain
vector sum, defined at every angle. At exactly antipodal orientation
it is $\mathbf{0}$ --- a computed answer, not a failure, and the same
absence of a favoured direction that \hyperref[subsec:axiom_null]{Axiom~2}
reports as black. Its magnitude is a free confidence signal, falling
smoothly from $2$ (identical directions) to $0$ (opposed).

\paragraph{Assembly never repels.}
For two units of nonzero strength considered alone, $U_{AB}<0$
always: there is no orientation, no pair of strengths, and no
separation at which isolated assembly pushes rather than pulls. This
is not an assumption but a property of
Equation~\eqref{eq:assembly_energy}, whose bracket cannot fall below
$2$. The one boundary is a black unit
(\hyperref[subsec:axiom_null]{Axiom~2}): at $r_A=0$ the energy is exactly
zero, no relation being formed with something absent, which is what
black already means elsewhere. This is taken as the definition,
$U_{AB}=0$ whenever $r_Ar_B=0$, because the formula alone gives $0/0$
for two black units on one seat. Whether a structure of three or more
units can hold members apart is a question for the enclosure, not
for the pairwise verb.

\paragraph{Polarity.}
The difference $\lvert r_A - r_B\rvert$ distinguishes a balanced bond
from a lopsided one, and so supplies a criterion for reading
$A\bowtie B$ as mutual or as actor-and-patient in the manner of $T$.
It is a label on the relation, not a difference in what the relation
does: Equation~\eqref{eq:assembly_energy} is symmetric under exchange
of $A$ and $B$, and $U_{AB}=U_{BA}$ exactly.

\paragraph{Where assembly sits among the verbs.}
Assembly is the first operator in this paper that requires anything
outside the 4-tuple. Gather, remove, Mix, rotation, and offset all
take units on the sphere and return units or states on the same
sphere; none of them asks \emph{where} a unit is, because nothing in
$(r,\hat n,\chi,\tau)$ answers that. Assembly asks, and so is stated
here as a verb but completed in Section~\ref{sec:enclosure}, which
supplies the seat $\vec P$ that $d_{AB}$ measures between. Read
without an enclosure, Equation~\eqref{eq:assembly_energy} still
gives a well-defined ordering of which pairs favour each other most
strongly at equal separation.

\subsection{Axiom 7: Informational vectors and derived per-atom metrics}
\label{subsec:axiom_v1_metrics}
The teaching atom of \hyperref[subsec:axiom1]{Axiom~1} is one 4-tuple
($N=1$). An atom may also be a \emph{bundle}: $N$ informational
vectors $v_i=(r_i,\hat n_i,\chi_i)$, each carrying its own
phase-color $\chi_i\in[0,2\pi)$ alongside $(r_i,\hat n_i)$ and
sharing the atom's own sort $\tau$ --- $v_i$ together with that
shared $\tau$ forming, per vector, the same 4-tuple
\hyperref[subsec:axiom1]{Axiom~1} states, still capped at $r_i\le
1$. $N$ is the atom's info
number. This is not a second species of unit and not a field on the
globe. It is several teaching atoms held under one closure ---
distinct phase-colors among them coexist without being forced
through Mix, the bundle's own instance of two identities occupying
one direction without overwriting one another
(\hyperref[subsec:axiom1]{Axiom~1}). At most one further,
\emph{derived} vector may sit alongside them in the same closure
unit:
\[
V_{\text{derived}} \;=\; \sum_{i=1}^{N} r_i\hat n_i,
\]
the raw, unnormalized sum of the informational vectors' geometric
part, computed after and from them, never including itself; $\chi_i$
does not enter this sum, the same way $\tau$ does not participate in
$T$'s rotation --- available wherever a per-vector reading is named,
silent in the seven metrics below, none of which is stated in terms
of it. $V_{\text{derived}}$ is
uncapped and excluded from $N$: it needs no $\tau$-justified
admissibility of its own (\hyperref[subsec:axiom0_admissibility]{Axiom~0}),
the same status the companion paper\footnote{\emph{Pattern
Arithmetic: Inference from Bowls and Paths}, which builds inference
machinery --- reference atoms, classification, this same
$\tau$-exempt status for uncapped witness-sums --- on top of the
atoms and operators defined here.} already grants its own uncapped
witness-sums. Its magnitude and direction read as an effective
$(r,\hat n)$ pair --- $r=\lvert V_{\text{derived}}\rvert$, $\hat n =
V_{\text{derived}}/r$ --- standing in for the atom wherever a weight
or a direction is required: Assembly's $r_Ar_B$
(\hyperref[subsec:axiom_assembly]{Axiom~6}), a camp body's alignment ---
neither involves a rotation or $\mathcal{S}$-membership. It may act
as the actor $w$ of $T$ (\S\ref{subsec:transform_definition}, whose
domain is stated as $w\in\mathcal{S}$) only after an explicit cast
back into $[0,1]$ at the call site --- the clip already used for
Mix's own uncapped sum (\hyperref[subsec:axiom_interference]{Axiom~5}), or,
with no new machinery at all, atom coherence itself, $\lvert
V_{\text{derived}}\rvert/S_{\text{atom}}$. Uncast, it is not a valid
$w$: three vectors at $r=0.5$, perfectly aligned, give $\lvert
V_{\text{derived}}\rvert=1.5$, outside $\mathcal{S}$ entirely, and
$\alpha(1.5)=270^\circ$ is not a rotation this axiom's $T$ takes ---
not because $\mathrm{SO}(3)$ forbids it, but because $w=1.5$ was
never $T$'s to accept. Until a cast is named at the call site, $T$
takes teaching atoms only; every worked $T$ path in this paper takes
one, and a bundle-as-actor example is not claimed here. For $N=1$,
$V_{\text{derived}}$ coincides with the atom's own vector, already
in $\mathcal{S}$, and nothing changes.

Seven scalar metrics follow from the same $N$ informational vectors,
each answering a different question, none used by this Axiom's own
Assembly (\S\ref{subsec:axiom_assembly}) or by Motion
(Section~\ref{sec:enclosure}) unless named:
\[
S_{\text{atom}} = \sum_{i=1}^N r_i, \qquad
\mathcal{V}_{\text{atom}} = \sum_{i=1}^N r_i^3, \qquad
\gamma_{\text{atom}} = \frac{\mathcal{V}_{\text{atom}}}{N}.
\]
$S_{\text{atom}}$, atom strength, is the correct denominator for
atom-level coherence, $\lvert V_{\text{derived}}\rvert /
S_{\text{atom}}$, structurally identical to the companion paper's own
mean resultant length one level down --- reading $1.0$ for perfectly
aligned informational vectors and falling smoothly toward $0$ as they
scatter. $\mathcal{V}_{\text{atom}}$, atom volume, is strictly
smaller than $S_{\text{atom}}$ whenever any $r_i<1$, a forced
consequence of the $r\le 1$ cap rather than an assumption: cubing a
number below $1$ always shrinks it, exactly as treating $r$ as a
literal radius and $r^3$ as the volume it encloses would predict.
$\gamma_{\text{atom}}$, mean vector volume, is the average per-vector
share of that total. None of the three participates in Motion; atom
volume, in particular, is deliberately excluded from it and exists
here purely as a descriptive metric of the atom. Earlier drafts
called $\mathcal{V}_{\text{atom}}$ ``mass,'' $\gamma_{\text{atom}}$
``density,'' and $S_{\text{atom}}$ ``atom energy'' --- all three
retired as misleading: the first two because neither participates in
momentum, kinetic energy, or any calculation the word ``mass'' would
ordinarily license, and the third because a genuine notion of energy
exists below and is not $S_{\text{atom}}$.

A fourth quantity is not derived from strength by squaring it, but
stands on its own, the way variance is defined directly rather than
as squared standard deviation:
\begin{equation}
    \mathcal{E}_{\text{atom}} \;=\; \sum_{i=1}^N r_i^2.
    \label{eq:information_energy}
\end{equation}
$\mathcal{E}_{\text{atom}}$ is the combined squared magnitude the
atom's informational vectors would carry if every pair were mutually
orthogonal --- confirmed by direct construction, computing
$\mathcal{E}_{\text{atom}}$ from a set of strengths matches $\lvert
V_{\text{derived}}\rvert^2$ computed from those same strengths placed
at mutually orthogonal directions, exactly. This is forced, not
chosen: an exact identity places $\mathcal{E}_{\text{atom}}$ inside
the derived vector's own squared magnitude,
\[
\lvert V_{\text{derived}}\rvert^2 \;=\; \mathcal{E}_{\text{atom}}
\;+\; 2\sum_{i<j} r_ir_j(\hat n_i\cdot\hat n_j),
\]
verified to six decimal places. The structure matches the variance
of a sum, $\operatorname{Var}(\sum X_i) = \sum\operatorname{Var}(X_i)
+ 2\sum_{i<j}\operatorname{Cov}(X_i,X_j)$, term for term: $r_i^2$
stands where $\operatorname{Var}(X_i)$ stands, each vector's own
self-contained contribution; $r_ir_j(\hat n_i\cdot\hat n_j)$ stands
where $\operatorname{Cov}(X_i,X_j)$ stands, how two vectors relate.
$\mathcal{E}_{\text{atom}}$ is the part that remains when every cross
term vanishes --- the atom's intrinsic energy, prior to any question
of how its own vectors align. This is distinct from the companion
paper's own bowl energy (a sum across many separate witnesses):
bowl energy sums across atoms, $\mathcal{E}_{\text{atom}}$ sums
squares within one.

A fifth metric answers a question the other four cannot. Atom
coherence above is a ratio of magnitudes and so cannot fall below
zero: it reads $0$ both for informational vectors that merely fail
to reinforce and for vectors that actively cancel. The
strength-weighted mean cosine similarity across all internal pairs
separates these:
\begin{equation}
    \alpha_{\text{atom}} \;=\;
    \frac{\sum_{i<j} r_i r_j\,(\hat n_i\cdot\hat n_j)}
         {\sum_{i<j} r_i r_j}.
    \label{eq:internal_alignment}
\end{equation}
Each $\hat n_i\cdot\hat n_j$ is already a cosine similarity, since
informational vectors are unit-directed; the denominator makes the
aggregate comparable across atoms of differing $N$ and differing
strengths, which the unnormalized sum is not --- two atoms with
comparably tight internal arrangement but $N=2$ and $N=5$ show how
little the raw sum alone says: $r=0.85$ throughout, pairwise
directions $5.7^\circ$ apart, gives a raw sum of $+0.72$ at $N=2$
(one pair) against $+7.14$ at $N=5$ (a five-fold symmetric cone of
the same half-angle, ten pairs) --- a tenfold spread driven purely by
pair count, not by any real difference in agreement, which
$\alpha_{\text{atom}}$ correctly reports as $0.995$ and $0.988$,
nearly identical. $\alpha_{\text{atom}}=+1$ exactly when all
informational vectors align; it is negative when they genuinely
oppose, reading $-1/3$ for two antipodal pairs on orthogonal axes
(four unit vectors, six pairs, two at cosine $-1$ and four at
cosine $0$). Its floor is not $-1$ in general: three
or more vectors in three dimensions cannot be mutually opposed, so
achievable disagreement tightens as $N$ grows while achievable
agreement does not --- the unnormalized sum has maximum
$\sum_{i<j}r_ir_j$, growing as $N^2$, against a hard minimum of
$-\tfrac{1}{2}\mathcal{E}_{\text{atom}}$, growing only as $N$, which
follows directly from $\lvert V_{\text{derived}}\rvert^2 =
\mathcal{E}_{\text{atom}} + 2\sum_{i<j} r_ir_j(\hat n_i\cdot\hat n_j)
\ge 0$. Agreement compounds; disagreement saturates.
$\alpha_{\text{atom}}$ is read alongside atom coherence, not in place
of it: the first asks whether the vectors agree, the second how much
of their strength survives into one resultant.

A sixth metric does for magnitude what $\alpha_{\text{atom}}$ does
for direction. $S_{\text{atom}}$ and $\mathcal{V}_{\text{atom}}$ both
report a total and are silent on how that total is spread: two atoms
can carry near-identical strength ($2.200$ against $2.150$ in a
checked case) while one holds it evenly across its vectors and the
other concentrates most of it in two. Writing $s_i = r_i^3 /
\mathcal{V}_{\text{atom}}$ for the share of volume held by vector
$i$:
\begin{equation}
    \sigma_{\text{atom}} \;=\;
    \frac{N^2}{N-1}\cdot\operatorname{Var}(s_1,\ldots,s_N),
    \qquad N\ge 2,
    \label{eq:volume_concentration}
\end{equation}
with $\sigma_{\text{atom}}=0$ by convention at $N=1$.
$\sigma_{\text{atom}}=0$ exactly when volume is spread evenly and $1$
when a single vector holds all of it, at every $N$ --- verified at
$N=2,3,4,6,8$ for both extremes. The normalization is what earns this
form over the simpler candidates: the largest vector's share and the
Herfindahl index were both tested and rejected, since each reads a
perfectly even atom as $1/N$ and so slides from $0.5$ to $0.125$ as
$N$ grows from $2$ to $8$, confounding evenness with size.
$\mathcal{V}_{\text{atom}}$ rather than $S_{\text{atom}}$ supplies
the shares deliberately: cubing amplifies concentration, so a single
dominant vector at $r=0.95$ against three at $r=0.2$ registers a
share of $0.973$, where linear strength would report the same atom
as unremarkable.

A seventh metric answers a question none of the first six can: not
how strength or volume is shared among distinct vectors, but whether
many vectors independently arrive at nearly the same magnitude.
$\sigma_{\text{atom}}$ is blind to this by construction --- four
hundred vectors at an identical $r=0.65$ each receive their own tiny
individual share, indistinguishable to
Equation~\eqref{eq:volume_concentration} from four hundred vectors
scattered across the full range; checked directly, a genuinely
clumped population registered as \emph{more} even
($\sigma_{\text{atom}}=0.00032$) than a uniform spread of the same
size ($0.00214$). Fixing this needs a shell width $w$, a named,
configurable parameter rather than a derived one, defaulting to
$w=0.1$ --- ten shells spanning $[0,1]$. Binning the $N$
informational vectors by $\lfloor r_i/w\rfloor$ into $K$ occupied
shells with populations $n_1,\ldots,n_K$ and $p_k = n_k/N$:
\begin{equation}
    \sigma_{\text{shell}}(w) \;=\;
    \begin{cases}
        1, & K=1 \\[4pt]
        \dfrac{K^2}{K-1}\cdot\operatorname{Var}(p_1,\ldots,p_K), & K\ge 2.
    \end{cases}
    \label{eq:shell_concentration}
\end{equation}
The $K=1$ case is stated separately because the general form has a
removable singularity there: every vector sharing one shell is
maximal concentration, not the zero the formula would otherwise
return. $\sigma_{\text{shell}}=0$ when occupied shells are evenly
populated, at every $K$ tested; $\sigma_{\text{shell}}=1$ when every
vector shares one shell. Checked at $w=0.1$ and $w=0.05$: well
separated shells (as in $10@0.98,\,3@0.3,\,400@0.65$, giving
$\sigma_{\text{shell}}=0.908$ at either width) are stable to the
choice, while genuinely uniform populations are not --- occupied
shell count nearly doubled between the two widths in a checked
uniform case, though the concentration score itself stayed
negligible either way. $w$ should be reported alongside
$\sigma_{\text{shell}}$ whenever it is used.


\section{Transformation Algebra: From Identity to Expression}
\label{sec:transformation_algebra}

\subsection{Motivation: Two Layers of Meaning}
\label{subsec:transform_motivation}
The operators established so far ($\oplus$, $\ominus$) answer the question
\emph{who is present?} Combining Sky, Sun, and Wind under $\oplus$
yields one order-independent bowl --- a bag of seats, not a blended
hue. Paint is Mix ($\boxplus$), a later reading. Analogies such as Sun-Day-Moon are \emph{not} gathers; they are deferred to the offset operator of Section~\ref{sec:analogical_offset}. Gather $\oplus$ does not, and structurally
cannot, distinguish a poem in which the sun rises through a wind-blown
sky from a poem in which wind extinguishes a darkening sun --- both
share the identical bowl
$R = \text{Sky} \oplus \text{Sun} \oplus \text{Wind}$, since $\oplus$
is commutative and idempotent by \hyperref[subsec:axiom_idempotence]{Axiom~4}.

The distinction between these two poems is not a difference in
\emph{ingredients}; it is a difference in \emph{how one state acts
upon another, and in what order}. This is a \emph{relational}
operation, and it requires a second operator, $T$, layered on top of
$\oplus$: where $\oplus$ builds nouns, $T$ builds sentences. A poem is
not the bowl $R$ itself, but a specific trajectory traced across
$\mathcal{S}$ by successive transformations applied to a register occupant.

\subsection{Formal Definition of the Transformation Operator}
\label{subsec:transform_definition}
Rather than introduce $T$ as an unrelated new primitive, it is defined
using the same rotational machinery already given in
\S\ref{subsec:axiom_rotation}: every state $w \in \mathcal{S}$
acts on another state $x \in \mathcal{S}$ by rotating $x$'s direction
about $w$'s own axis, by an angle proportional to $w$'s strength, while
additively shifting $x$'s phase by $w$'s phase:
\begin{equation}
    T_w(x)
    \;=\;
    \Bigl(
        r_x,\;\;
        \mathrm{Rot}_{\hat n_w}\!\bigl(\alpha(r_w)\bigr)\hat n_x,\;\;
        (\chi_x + \chi_w) \bmod 2\pi,\;\;
        \tau_x
    \Bigr)
    \label{eq:transform_def}
\end{equation}
where $\hat n_w$ is the unit axis of $w$,
$\mathrm{Rot}_{\hat n}(\alpha) \in \mathrm{SO}(3)$ is the
rotation of \S\ref{subsec:axiom_rotation} about that axis by angle
$\alpha$, and $\alpha(r_w) = \pi r_w$ is the simplest monotonic choice
mapping strength to rotation angle (a full-strength unit, $r_w = 1$,
rotates by the maximal $180^\circ$). Magnitude $r_x$ and sort $\tau_x$ are left invariant:
a transformation reorients and recolors a state, it does not
strengthen or weaken it, and it does not recast it --- those remain
the roles of Mix (strength) and of an explicit cast (sort).

\subsection{Derived Properties}
\label{subsec:transform_properties}
The following properties follow from Equation~\eqref{eq:transform_def}
rather than being separately postulated.

\paragraph{Identity.} If $w = \mathbf{0}$ (the Null State of
\hyperref[subsec:axiom_null]{Axiom~2}), then $\alpha(0) = 0$ and the phase
shift is $0$, so $T_{\mathbf{0}} = \mathrm{id}$.

\paragraph{Invertibility.}
The inverse of $T_w$ is not ``the same axis and the same $+\alpha$ again.'' Same axis, same positive angle, repeats the turn. Inverse is $-\alpha$ on the same axis, or $+\alpha$ on the antipodal axis, together with $-\chi_w$. Define
\begin{equation}
    w^{*} = \bigl(r_w,\,-\hat n_w,\,-\chi_w \bmod 2\pi,\,\tau_w\bigr),
\end{equation}
where $-\hat n_w$ is the antipode $(\pi-\theta_w,\,\phi_w+\pi)$.
Then $T_{w^{*}}\!\big(T_w(x)\big) = x$ for all $x$. As rotations, only $\alpha=0$ ($r_w=0$) and $\alpha=180^\circ$ ($r_w=1$) equal their own reverse; the inverse state $w^{*}$ is still required so that phase is cancelled too. This undoes a transformation on a state that persists. It is distinct from $\{A\}\ominus A=\emptyset$ (last shell removed, leaving the empty bowl) and from Mix of equal anti-phase marks (a reading into black). Neither of those is an inverse of $T$.

\paragraph{Associativity.} Since each $T_w$ is a function
$\mathcal{S} \to \mathcal{S}$, composition of transformations is
associative automatically, by the associativity of function
composition; no additional axiom is required.

\paragraph{Non-commutativity.} Rotations about distinct axes in
$\mathrm{SO}(3)$ generically do not commute: for $a, b$ with distinct
directions $(\theta,\phi)$,
\begin{equation}
    T_a\big(T_b(x)\big) \;\neq\; T_b\big(T_a(x)\big).
    \label{eq:noncommute}
\end{equation}
This is not an additional assumption; it is inherited directly from
the well-established non-commutativity of $\mathrm{SO}(3)$, already
invoked by \S\ref{subsec:axiom_rotation}. It is this property,
and only this property, that allows two poems built from the same
bowl to diverge: order of application is meaningful even
though the underlying identity is not.

\subsection{Beyond three dimensions}
\label{subsec:transform_higher_dim}
Equation~\eqref{eq:transform_def} needs only $\hat n_w$ to fix a
rotation because $\mathrm{SO}(3)$'s generators happen to form a
$3$-dimensional space, exactly matching $\mathbb{R}^3$ --- not
because one vector generally suffices to specify a rotation. It does
not: a \emph{simple} rotation, one that fixes as much as possible and
turns a single remaining plane, fixes a subspace of dimension
$N_{\mathrm{dim}}-2$ and rotates its $2$-dimensional orthogonal
complement. The ladder climbs one dimension at a time. A circle
turns about a fixed centre; stacking that centre through the circular
slices of a sphere gives the sphere's $1$-dimensional axis; stacking
again gives the $2$-dimensional axis plane of the sphere in four
dimensions, and again the $3$-dimensional fixed space of the sphere
in five:
\begin{center}
\begin{tabular}{llll}
Space & Sphere & Fixed part (the ``axis'') & What turns \\
\hline
$\mathbb{R}^2$ & $S^1$, a circle & a point, dimension $0$ & one plane \\
$\mathbb{R}^3$ & $S^2$ & a line, dimension $1$ & one plane \\
$\mathbb{R}^4$ & $S^3$ & a plane, dimension $2$ & one plane \\
$\mathbb{R}^5$ & $S^4$ & a $3$-D space, dimension $3$ & one plane \\
$\mathbb{R}^{N_{\mathrm{dim}}}$ & $S^{N_{\mathrm{dim}}-1}$ & dimension $N_{\mathrm{dim}}-2$ & one plane \\
\end{tabular}
\end{center}
At every rung the fixed part grows by one dimension and the turning
part stays a single plane. (A sphere is named by its own dimension,
not by the space around it: the sphere in four dimensions is $S^3$.)
Three dimensions is not where axes stop needing more than a point ---
it is where they still fit in one vector.

For $\mathcal{S}=S^{N_{\mathrm{dim}}-1}\subset\mathbb{R}^{N_{\mathrm{dim}}}$,
$w$ supplies the fixed subspace $F$ directly, as $N_{\mathrm{dim}}-2$
mutually orthogonal informational vectors from its own bundle
(\hyperref[subsec:axiom_v1_metrics]{Axiom~7}) rather than as one $\hat n_w$:
\begin{equation}
    T_w(x)
    \;=\;
    \Bigl(
        r_x,\;\;
        \hat n_x^F + \mathrm{Rot}_{\alpha(r_w)}\!\bigl(\hat n_x^P\bigr),\;\;
        (\chi_x + \chi_w) \bmod 2\pi,\;\;
        \tau_x
    \Bigr)
    \label{eq:transform_def_general}
\end{equation}
where $\hat n_x^F$ is $\hat n_x$'s component in $F$, left untouched,
and $\hat n_x^P$ its component in the remaining $2$-dimensional plane
$F^\perp$, rotated within that plane by $\alpha(r_w)=\pi r_w$ exactly
as before. This requires $N_{\mathrm{dim}}-2$ mutually orthogonal
arrows in the actor's bundle, so $N\ge N_{\mathrm{dim}}-2$ is
necessary but not by itself sufficient: an atom acting as
$w$ in more than three dimensions needs a bundle rich enough to
supply its own axis, not just a direction. At $N_{\mathrm{dim}}=3$,
$F$ is one vector and Equation~\eqref{eq:transform_def_general} is
Equation~\eqref{eq:transform_def} exactly --- checked directly, not
merely argued: the same inputs give the same output to machine
precision.

\begin{center}
\begin{tabular}{lc}
Stage & Orthonormal arrows required on the \emph{actor} \\
\hline
$S^2$ (3-D) & $1$ --- the teaching atom, unchanged \\
$S^3$ (4-D) & $2$ \\
$S^{11}$ (12-D) & $10$ \\
\end{tabular}
\end{center}
A teaching atom, $N=1$, is therefore a legal actor only at
$N_{\mathrm{dim}}=3$; this is why three dimensions is the standard
this paper keeps, not an arbitrary default --- a single arrow can
walk without a bundle only there.

Checked beyond the coincidence of matching 3D: invertible by
$-\alpha$ on the same $w$ at every angle tested from $0$ to
$359^\circ$, not only small ones --- the naive alternative, borrowing
the missing axis data from $x$ instead of $w$, is invertible only
while $\alpha$ stays short of the angle between $w$ and $x$, and
silently fails past it, since the borrowed direction depends on
which side of $w$ the rotation has already passed. Sourcing $F$ from
$w$ alone avoids this: the fixed subspace never depends on what is
being acted on, the way $\mathrm{SO}(3)$'s axis never did.
Non-commutativity for distinct actors and angle-composition for
repeated ones both continue to hold, checked directly at
$N_{\mathrm{dim}}=4$, $5$, $8$, $12$, $20$, and $50$ on a
configuration that needs no seed. With $e_i$ the coordinate axes, actor $w_1$ fixes
$e_1,\ldots,e_{N_{\mathrm{dim}}-2}$ and so turns the plane of
$e_{N_{\mathrm{dim}}-1},e_{N_{\mathrm{dim}}}$; actor $w_2$ fixes
$e_2,\ldots,e_{N_{\mathrm{dim}}-1}$ and so turns the plane of
$e_1,e_{N_{\mathrm{dim}}}$; each plane is oriented from its
lower-indexed axis to its higher. The occupant is
$\hat n_x=(1,2,\ldots,N_{\mathrm{dim}})$ normalized. Invertibility is
tested at $30^\circ$, $45^\circ$, $90^\circ$, $135^\circ$,
$180^\circ$, $200^\circ$, $270^\circ$, and $359^\circ$; the order
gap is $\lvert T_{w_2}T_{w_1}\hat n_x-T_{w_1}T_{w_2}\hat n_x\rvert$
with both actors at $90^\circ$:
\begin{center}
\begin{tabular}{ccc}
$N_{\mathrm{dim}}$ & largest inversion error & order gap \\
\hline
$4$ & $1.1\times10^{-16}$ & $0.683$ \\
$5$ & $1.1\times10^{-16}$ & $0.688$ \\
$8$ & $1.1\times10^{-16}$ & $0.649$ \\
$12$ & $5.6\times10^{-17}$ & $0.584$ \\
$20$ & $1.1\times10^{-16}$ & $0.489$ \\
$50$ & $5.6\times10^{-17}$ & $0.331$ \\
\end{tabular}
\end{center}
Inversion is exact to rounding at every angle tested, and the two
orders disagree at every dimension tested. Both results stay on the unit sphere.

One choice is left open, and is a convention rather than a gap:
which orthonormal basis of $F^\perp$ orients the rotation (which
direction counts as $+\alpha$) is fixed by $w$'s own
$N_{\mathrm{dim}}-2$ vectors, but the specific construction is not
unique --- no more than the right-hand rule is the only possible
convention for $\mathrm{SO}(3)$'s own cross product. Any consistent
choice, stated once, suffices; \S\ref{subsec:transform_example_4d}
states and exercises one.

\subsection{Poems as Transformation Trajectories}
\label{subsec:transform_trajectories}
A poem $\pi$ built from a register occupant $x$ and an ordered sequence of
transforming states $w_1, w_2, \ldots, w_n$ (typically drawn from the bowl) is the trajectory
\begin{equation}
    \pi = \big(x;\, w_1, \ldots, w_n\big): \quad
    x \;\longrightarrow\; T_{w_1}(x) \;\longrightarrow\; T_{w_2}T_{w_1}(x) \;\longrightarrow\; \cdots \;\longrightarrow\; T_{w_n}\cdots T_{w_1}(x).
    \label{eq:poem_trajectory}
\end{equation}
Two poems may share an identical bowl (the same theme, the
same ``ingredients'') while tracing entirely different trajectories
--- this is the formal sense in which many poems arise from the same
units of truth. Conversely, because rotation composition is
many-to-one, two distinct trajectories may terminate at the same
endpoint; a poem's identity is therefore carried by its full path
through $\mathcal{S}$, not by its final state alone.

Unlike $\oplus$, $T$ does not collapse a repeated actor: $T_w\circ
T_w\neq T_w$ in general, since composing a rotation with itself
composes the angle, $T_w\circ T_w = T_{w'}$ for some $w'\neq w$ rather
than returning $T_w$. A path that visits the same actor $n$ times is
therefore $n$ genuine steps, not one step counted $n$ times, and its
length recovers that count without any coordinate carrying it. This is
the sense in which the path, unlike the bowl, does not discard
multiplicity (Section~\ref{subsec:axiom_idempotence}).

A step of the trajectory may carry more than the bare actor:
$\pi=\bigl((x;w_1,m_1),\ldots,(w_n,m_n)\bigr)$ for some label $m_i$ ---
a discrete time bucket, say. This changes nothing about the atom: $m_i$
is metadata on the walk, not a coordinate of any $x$ or $w_i$, and
Equation~\eqref{eq:poem_trajectory} runs exactly as before on the
$w_i$ alone. What such a label buys, and does not buy, follows
directly from the discreteness point above: a handful of buckets
behaves like a sort tag and does not block $A\oplus A=A$ where it
applies, but a bucket is a fact about how a particular walk was
assembled, not a property its occupant carries into the next one.

\subsection{Worked Numerical Example: Sun, Wind, and Sky}
\label{subsec:transform_example}
Let Sky be the state acted upon, with direction
$(\theta,\phi) = (45^\circ, 45^\circ)$, i.e.\ Cartesian axis
$(0.500, 0.500, 0.707)$. Let Sun define an axis along
$(\theta,\phi) = (90^\circ, 0^\circ)$, i.e.\ $(1,0,0)$, with strength
$r_{\text{Sun}} = 1.0$, giving rotation angle $\alpha_{\text{Sun}} =
180^\circ$. Let Wind define an axis along $(\theta,\phi) = (90^\circ,
90^\circ)$, i.e.\ $(0,1,0)$, with strength $r_{\text{Wind}} = 0.7$,
giving $\alpha_{\text{Wind}} = 126^\circ$.

\emph{Order 1 --- Sun acts first, then Wind} (the sun rising through a
wind-blown sky):
\begin{equation}
    T_{\text{Wind}}\big(T_{\text{Sun}}(\text{Sky})\big) \;\approx\; (-0.866,\, -0.500,\, 0.011)
\end{equation}

\emph{Order 2 --- Wind acts first, then Sun} (wind extinguishing a
darkening sun):
\begin{equation}
    T_{\text{Sun}}\big(T_{\text{Wind}}(\text{Sky})\big) \;\approx\; (0.278,\, -0.500,\, 0.820)
\end{equation}

The two endpoints are clearly distinct, confirming
Equation~\eqref{eq:noncommute}: identical primitives, identical
bowl, opposite structural readings. Notably, if $\chi_{\text{Sun}}
= 0^\circ$ and $\chi_{\text{Wind}} = 90^\circ$, the accumulated phase
shift is $270^\circ$ in \emph{both} orderings, since phase addition is
commutative. The two poems therefore share the same blended color/mood
even as their spatial --- structural --- meaning diverges. This
asymmetry (shared affect, divergent structure) is a numerical check of
the algebra, not a theory of verse.

\subsection{Worked Numerical Example: A Four-Dimensional Actor}
\label{subsec:transform_example_4d}
The construction of \S\ref{subsec:transform_higher_dim} is checked
here by hand, not only by claim. Let $N_{\mathrm{dim}}=4$. An actor
$w$ supplies its required $N_{\mathrm{dim}}-2=2$ orthonormal arrows as
the first two standard basis vectors, $e_1=(1,0,0,0)$ and
$e_2=(0,1,0,0)$, fixing $F=\mathrm{span}(e_1,e_2)$ and leaving the
active plane $F^\perp=\mathrm{span}(e_3,e_4)$. Take $r_w=0.5$, so
$\alpha(r_w)=90^\circ$. Let the occupant be
\[
\hat n_x = \Bigl(\tfrac12,\ \tfrac12,\ \tfrac{1}{\sqrt2},\ 0\Bigr),
\]
a unit vector with a nonzero component in both $F$ and $F^\perp$, so
the example shows both halves of Equation~\eqref{eq:transform_def_general}
at once. Its $F$-component, $(\tfrac12,\tfrac12,0,0)$, is exactly
what $T_w$ leaves untouched; its $F^\perp$-component,
$(0,0,\tfrac{1}{\sqrt2},0)$, is what turns by $90^\circ$ within
$\mathrm{span}(e_3,e_4)$, landing on $(0,0,0,\tfrac{1}{\sqrt2})$.
Summing:
\begin{equation}
    T_w(\hat n_x) \;=\; \Bigl(\tfrac12,\ \tfrac12,\ 0,\ \tfrac{1}{\sqrt2}\Bigr).
\end{equation}
Applying $T_w$ again with $-\alpha$ recovers the $F^\perp$-component
exactly, $(0,0,\tfrac{1}{\sqrt2},0)$, since $F$ was never touched to
begin with:
\begin{equation}
    T_{w^*}\bigl(T_w(\hat n_x)\bigr) \;=\; \Bigl(\tfrac12,\ \tfrac12,\ \tfrac{1}{\sqrt2},\ 0\Bigr) \;=\; \hat n_x,
\end{equation}
confirming invertibility at a genuine right angle, not only for the
small perturbations a reader might otherwise suspect were doing the
work. The orientation convention of
\S\ref{subsec:transform_higher_dim} is exercised here for the first
time: $f_1=e_3$, $f_2=e_4$, in that order, is the specific choice
this example makes, and a different ordering would rotate the other
way.

\subsection{Scope}
\label{subsec:transform_scope}
$T$ is deliberately left unconstrained beyond
Equation~\eqref{eq:transform_def}: any sequence of any transforming
states is a formally valid trajectory in three dimensions, or in
$N_{\mathrm{dim}}$ dimensions provided each actor's own bundle is
rich enough to name its $F$ (\S\ref{subsec:transform_higher_dim}); an
underqualified actor is not a shorter path, it is not a legal $w$ at
all. Pattern Arithmetic in this form
supplies the substrate on which expression is built --- the
combinatorial space of possible trajectories --- and does not yet
supply a grammar constraining which trajectories are coherent,
well-formed, or meaningful to a human reader. Whether such constraints
belong inside this framework (as further axioms on admissible
sequences of $T$) or outside it (as a learned or externally specified
grammar layered on top) is left open.

\section{Analogical Arithmetic: The Offset Operator}
\label{sec:analogical_offset}

Offset is another verb. It is not $\oplus$ (no gather), not Mix
(no phasor blend), and not $T$
(no spin of an actor on a path). It copies a \emph{walk} from one
seat to another. The landing is a new direction on the same globe.
A dictionary may later pin a name on that landing. The walk does
not mint the name.

\subsection{Recipe: logarithm, carry, exponential}
\label{subsec:offset_logexp}
Write directions as unit vectors $\hat n\in S^2\subset\mathbb{R}^3$.
The walk from $p$ to $q$ is the tangent vector
\begin{equation}
    \theta = \arccos(\hat p\cdot\hat q),
    \qquad
    \log_{\hat p}(\hat q)
    = \frac{\theta}{\sin\theta}\bigl(\hat q-(\hat p\cdot\hat q)\,\hat p\bigr)
    \quad (\theta\neq 0),
    \label{eq:sphere_log}
\end{equation}
with $\log_{\hat p}(\hat p)=\mathbf{0}$. Its length is the geodesic
angle. The reverse map, from a tangent vector $v$ at $p$, is
\begin{equation}
    \exp_{\hat p}(v)
    = \cos\|v\|\,\hat p + \sin\|v\|\,\frac{v}{\|v\|}
    \quad (v\neq\mathbf{0}),
    \label{eq:sphere_exp}
\end{equation}
and $\exp_{\hat p}(\mathbf{0})=\hat p$. Given three seats
$m,k,w$ (source, source-image, new start), set
\begin{equation}
    v = \log_{\hat n_m}(\hat n_k),
    \qquad
    \hat n_{*}
    = \exp_{\hat n_w}\bigl(\mathrm{transport}_{m\to w}(v)\bigr).
    \label{eq:analogical_offset}
\end{equation}
The carry rotates the tangent vector by the geodesic angle about the
axis of $m\to w$:
\begin{equation}
    \theta_{mw} = \arccos(\hat n_m\cdot\hat n_w),
    \qquad
    \hat u
    = \frac{\hat n_m\times\hat n_w}{\|\hat n_m\times\hat n_w\|}
    \quad (\hat n_m\neq\pm\hat n_w),
    \label{eq:transport_axis}
\end{equation}
\begin{equation}
    \mathrm{transport}_{m\to w}(v)
    = \mathrm{Rot}_{\hat u}(\theta_{mw})\, v
    = v\cos\theta_{mw}
      + (\hat u\times v)\sin\theta_{mw}
      + \hat u\,(\hat u\cdot v)\,(1-\cos\theta_{mw}).
    \label{eq:transport}
\end{equation}
The axis is undefined when $\hat n_m=\pm\hat n_w$.
Coincident seats ($\hat n_m=\hat n_w$): the carry is $v$ itself.
Antipodal seats ($\hat n_m=-\hat n_w$): no unique short geodesic;
that case is a separate convention and is not used below.
For a first recipe, if $m$ and $w$ are not far, transport may be
omitted and $v$ used at $w$ as written. That case is crude and
marked as such. The full recipe is~\eqref{eq:transport}, so the
heading stays a ground heading.

Strength and colour are not invented by the walk. This paper takes
\begin{equation}
    r_{*}=r_w,
    \qquad
    \chi_{*}=(\chi_w+\chi_k-\chi_m)\bmod 2\pi,
    \qquad
    \tau_{*}=\tau_w.
    \label{eq:offset_rchi}
\end{equation}
The target keeps its $r$ and its sort. The mood increment of the pair
$m\to k$ is reused. Nothing here uses $\alpha=\pi r$ from $T$.
With~\eqref{eq:analogical_offset} and~\eqref{eq:offset_rchi} together,
offset returns a complete state, not a bare direction:
\[
\text{offset}:\;\mathcal{S}\times\mathcal{S}\times\mathcal{S}
\longrightarrow\mathcal{S},
\qquad
(m,k,w)\mapsto
\bigl(r_w,\;\hat n_{*},\;\chi_{*},\;\tau_w\bigr).
\]
It is therefore a verb on the same footing as $\oplus$, $\boxplus$ and
$T$: a landing may be gathered into a bowl, walked from, or used as
the start seat of a further offset. The question of \emph{whether
anything is seated there} is separate, and is settled by a dictionary,
not by the walk. $\tau_{*}=\tau_w$ because the question asked is what
stands to $w$ as $k$ stands to $m$; the answer is of $w$'s kind. A
walk copied between sorts would need an explicit cast
(Section~\ref{sec:encodings_sorts}), which offset does not perform.
Two seated units keep two strengths. If $\mathrm{Day}$ has $r=0.9$
and $\mathrm{Night}$ has $r=0.8$, both sit; there is no
$0.9+0.8=1.7$. That sum is ambient $\mathbb{R}^3$ arithmetic and
leaves the ball. $\oplus$ gathers the two seats. A later path may
turn a dial or clip a mix. It does not get to smash two $r$'s into
one illegal length as the price of seating them.

\paragraph{Beyond three dimensions.}
Equation~\eqref{eq:transport} is written with a cross product, which
exists only in three dimensions, but what it performs is the simple
rotation of \S\ref{subsec:transform_higher_dim}: it turns the plane
spanned by $\hat n_m$ and $\hat n_w$ through $\theta_{mw}$, carrying
$m$ onto $w$, and fixes every direction orthogonal to that plane.
Written that way it holds on any $S^{N_{\mathrm{dim}}-1}$:
\begin{equation}
    \begin{aligned}
    \mathrm{transport}_{m\to w}(v)
    = v &+ (\cos\theta_{mw}-1)\bigl[(v\cdot\hat n_m)\hat n_m + (v\cdot\hat b)\hat b\bigr]\\
        &+ \sin\theta_{mw}\bigl[(v\cdot\hat n_m)\hat b - (v\cdot\hat b)\hat n_m\bigr],
    \end{aligned}
    \label{eq:transport_general}
\end{equation}
where $\hat b = \bigl(\hat n_w - (\hat n_m\cdot\hat n_w)\hat n_m\bigr)/
\lVert \hat n_w - (\hat n_m\cdot\hat n_w)\hat n_m\rVert$ is the unit
direction of $w$ as seen from $m$ within their plane.
Unlike $T$, offset needs no bundle for this: it is handed two seats,
and the two seats name the turning plane themselves. At
$N_{\mathrm{dim}}=3$, Equation~\eqref{eq:transport_general} is
Equation~\eqref{eq:transport} exactly, and the equator example below
lands on the same point. In four dimensions, with
$\hat n_m=(1,0,0,0)$, $\hat n_k=(\tfrac{\sqrt3}{2},0,0,\tfrac12)$, and
$\hat n_w=(0,1,0,0)$, the step $v=\log_{\hat n_m}(\hat n_k)
=(0,0,0,\tfrac{\pi}{6})$ lies outside the turning plane, is carried
unchanged, and lands at $\hat n_*=(0,\tfrac{\sqrt3}{2},0,\tfrac12)$.

\subsection{A check on one equator, not a definition}
\label{subsec:offset_equator}
Place three arrows in the $xy$-plane:
\[
\hat n_m=(1,0,0),\quad
\hat n_k=\bigl(\tfrac{\sqrt{3}}{2},\tfrac12,0\bigr),\quad
\hat n_w=(0,1,0)
\]
($0^\circ$, $30^\circ$, $90^\circ$). Then
$\log_{\hat n_m}(\hat n_k)=(0,\pi/6,0)$.
Here $\theta_{mw}=90^\circ$ and $\hat u=(0,0,1)$, so
\eqref{eq:transport} is a quarter-turn about $\hat z$ and gives
$(-\pi/6,0,0)$. The exponential at
$\hat n_w$ lands at
\[
\hat n_{*}=\bigl(-\tfrac12,\,\tfrac{\sqrt{3}}{2},\,0\bigr)
\quad (120^\circ).
\]
On this slice the walk \emph{looks like} $90^\circ+30^\circ=120^\circ$.
That sum is a protractor check. It is not the operator, and it is
not a night. If a lexicon writes the names
$\mathrm{Sun},\mathrm{Day},\mathrm{Moon}$ on the first three seats,
it may write $\mathrm{Night}$ on the landing: whatever $\mathrm{Sun}$
is to $\mathrm{Day}$, $\mathrm{Moon}$ is to that fourth place. The
geometry only promised a fourth pin.
The same moral as weather: $\mathrm{Day}$ and $\mathrm{Night}$
are dictionary pins, not sums of luminaries. The step that carries
$\mathrm{Sun}$ to $\mathrm{Day}$ is a walk, and the walk is what is
copied; $\oplus$ of those seats does not land on a time of day.
Whether the two halves really are the same walk depends entirely on
how a lexicon seated them, which is a dictionary claim and not one
this paper makes. Offset only copies a walk; a later comparison
may find the landing \emph{near} an already seated name. It does
not mint the name.

\subsection{Offset is not $T$}
\label{subsec:offset_not_T}
$T_w(x)$ spins an occupant about an axis. Offset plants a new seat
from three seats. A path does not need~\eqref{eq:analogical_offset}.
A dictionary act does not need~\eqref{eq:transform_def}. Do not
encode the walk as a ghost actor whose sole job is to trigger $T$.

The popular ambient sum $\hat n_k-\hat n_m+\hat n_w$ then
normalize (Recipe~1) is the embedding fluke written as geometry.
It is recorded as a cheaper cousin, not as a derivation of the
landing. This paper uses Recipe~2 so the copied step
stays on the skin of the ball.

\section{Encodings, Sorts, and Mixed Documents}
\label{sec:encodings_sorts}

The unit sphere is a single bounded \emph{register}.
Numbers, words, and colors are not three different globes; they are
three encodings of a token onto the same globe. A scientific paragraph
already mixes those alphabets on one page. Pattern Arithmetic may do
the same, provided each loaded state carries a \emph{sort} and the
operators respect that sort.

\subsection{Two primary sorts}
\label{subsec:sorts}
A state $x\in\mathcal{S}$ is the typed atom of
\hyperref[subsec:axiom1]{Axiom~1}. What the geometric 4-tuple
$(r,\theta,\phi,\chi)$ \emph{means} depends on the sort $\tau$ and the
encoding used to write it:
\begin{itemize}
    \item \textbf{Numeric sort ($\mathtt{num}$).}
          The state represents a scalar. One toy encoding --- the only
          numeric encoding used in this paper --- stores an
          integer $n$ on a marked clock as
          $\chi = 2\pi n/N$ (arithmetic modulo $N$), at a fixed
          direction, with $r=1$.
          Another, stereo-style, is recorded as an option: it would
          send $x\to\pm\infty$ toward the
          poles so that unbounded reals still occupy the globe. It is
          not used below.
    \item \textbf{Lexical sort ($\mathtt{word}$).}
          The state is a labeled point in a toy dictionary living on
          the same globe. Similarity is the angle between points
          (cosine similarity). A bag of words is formed by $\oplus$.
          A sentence or poem is formed by $T$.
\end{itemize}
Schoolbook addition lives only in $\mathtt{num}$ mode.
Mix as a \emph{domain} reading lives in the Tier-1 sorts of
\hyperref[subsec:axiom0_admissibility]{Axiom~0}: $\mathtt{color}$ under this
paper's phasor inclusion, and $\mathtt{note\text{-}pitch}$, whose
octave equivalence makes the circle native rather than borrowed.
The two sorts worked below are $\mathtt{num}$ and $\mathtt{word}$;
neither takes a domain Mix.
Two $\mathtt{word}$ states may be mixed as geometry; that is not a
$\chi$-reading of language. The two equations may share a sphere.
They may not share an operator.

\subsection{Numeric mode: $7+5=12$ on the clock, and why $\oplus$ is not $12$}
\label{subsec:clock_encoding}
Fix the modulus $M=24$ and the encoding
\begin{equation}
    E_{\mathtt{num}}(k)
    \;=\;
    \bigl(r=1,\;\hat n=(1,0,0),\;\chi=2\pi k/24,\;\tau=\mathtt{num}\bigr),
    \qquad
    k\in\mathbb{Z}/24\mathbb{Z}.
    \label{eq:clock_encoding}
\end{equation}
Decode is the inverse chart: $k = M\chi/(2\pi)$, an integer hour mark.
Then $7$ sits at $\chi=105^\circ$, $5$ at $\chi=75^\circ$, and $12$ at
$\chi=180^\circ$. Schoolbook $+$ is decoded arithmetic:
\begin{equation}
    7+5
    \;=\;
    E_{\mathtt{num}}^{-1}\bigl(E_{\mathtt{num}}(7)\bigr)
    +
    E_{\mathtt{num}}^{-1}\bigl(E_{\mathtt{num}}(5)\bigr)
    \;=\; 12,
    \qquad
    E_{\mathtt{num}}(12)
    \;=\;
    (1,\,(1,0,0),\,180^\circ,\,\mathtt{num}).
\end{equation}
That is $7+5=12$ as a calculation on $\mathcal{S}$: decode, add in
$\mathbb{Z}$, encode. Gather is legal on the same pair but answers a
different question: $7\oplus 5=\{7,5\}$ records which numerals are
present and is not $12$. If one instead applies Mix
(Equation~\eqref{eq:interference_formal}) to the same two geometric
marks, the phasors are
\begin{equation}
    z
    \;=\;
    e^{i 7\pi/12} + e^{i 5\pi/12}
    \;=\;
    2\cos(\pi/12)\,e^{i\pi/2},
\end{equation}
so $r'=\min\bigl(1,\,2\cos 15^\circ\bigr)=1$ and $\chi'=90^\circ$.
The clock would read that mark as hour $6$, not hour $12$.
A bowl $\{\,E_{\mathtt{num}}(7),\,E_{\mathtt{num}}(5)\,\}$ is two seats,
still not $12$. This single contrast is why the type table exists:
$+$ answers \emph{how many}; $\oplus$ and Mix do not.

\subsection{A document is a sequence, not a blend}
\label{subsec:document_sequence}
Words and numbers coexist in reports because a document is an ordered
tape of tokens, not one mixed arrow. Write a paragraph as a sequence
of typed atoms
\begin{equation}
    (x_1,\ldots,x_m),
    \qquad
    x_i=(r_i,\,\hat n_i,\,\chi_i,\,\tau_i),
    \qquad
    \tau_i\in\{\mathtt{num},\mathtt{word}\},
    \label{eq:mixed_document}
\end{equation}
The sphere holds one current occupant of the register; the paragraph is the history of
loads and actions. A poem
(Section~\ref{subsec:transform_trajectories}) is one species of that
history: a $T$-walk on one occupant. Juxtaposition is another. Mix is
not.

They may meet in only three honest ways.

\paragraph{Juxtaposition.}
The default. Tokens sit in order and no algebra is applied between
them. In ``the sample reached $37.4^\circ\mathrm{C}$ at dawn,'' the
number $37.4$ is not blended with the word $\mathrm{dawn}$.

\paragraph{Typed action.}
The operator $T_w(x)$ is allowed across sorts only with an explicit
output sort. Two useful conventions, which the rest of this paper
adopts, are:
\begin{itemize}
    \item a word acting on a number yields a number-in-context
          (still $\mathtt{num}$);
    \item a number acting on a word yields a refined word
          (still $\mathtt{word}$).
\end{itemize}
A number may also set a numeric dial of $T$ --- for example the
rotation angle $\alpha(r)$ --- because that dial is already a scalar.
That is how ``wind of $12\,\mathrm{m/s}$'' can strengthen a turn
without becoming a synonym of $\mathrm{wind}$.

\paragraph{Cast.}
Rare and named. The map $E_{\mathtt{num}\to\mathtt{word}}(2)=\texttt{two}$
turns a digit into a spoken numeral; the inverse is parsing. Casts are
dictionary maps, not $\oplus$.

The illegal meeting is an untyped blend
\begin{equation}
    \mathrm{dawn} \oplus 37.4,
\end{equation}
which has no meaning until one side is first cast into the other's
sort.

\subsection{Operator legality}
\label{subsec:type_table}
Table~\ref{tab:type_table} records which buttons accept which pairs.
Here $+$ stands for schoolbook arithmetic on the numeric encoding
(decode, operate in $\mathbb{Z}$, encode; $+$ is the worked case of
Section~\ref{subsec:clock_encoding}), $\oplus$ is
idempotent gather, $\boxplus$ is Mix, $T$ is typed action, and $x-a+b$ is the analogical
offset of Section~\ref{sec:analogical_offset}.

\begin{table}[h]
\centering
\begin{tabular}{lccccc}
\hline
Pair & $+$ & $\oplus$ & $\boxplus$ & $T_w(x)$ & $x-a+b$ \\
\hline
$(\mathtt{word},\mathtt{word})$ & no & yes & yes$^{\dagger}$ & yes & yes \\
$(\mathtt{num},\mathtt{num})$   & yes & yes$^{\ddagger}$ & no  & no$^{*}$ & no$^{*}$ \\
$(\mathtt{word},\mathtt{num})$  & no & yes$^{\ddagger}$ & no  & typed & no \\
$(\mathtt{num},\mathtt{word})$  & no & yes$^{\ddagger}$ & no  & typed & no \\
$(\mathtt{color},\mathtt{color})$ & no & yes & yes & yes & open \\
$(\mathtt{note\text{-}pitch},\mathtt{note\text{-}pitch})$ & no & yes & yes & yes & open \\
$(\mathtt{chladni},\mathtt{chladni})$ & no & yes & yes$^{\S}$ & open & open \\
$(\tau_1,\tau_2)$ any other & no & yes$^{\ddagger}$ & open & open & open \\
\hline
\end{tabular}
\caption{Sort discipline on a single sphere. Every $\boxplus$ entry is
additionally subject to the place restriction of
\hyperref[subsec:axiom_interference]{Axiom~5}: Mix combines $\chi$ at a common
$\hat n$, and Mix across differing directions is open
(Section~\ref{subsec:mix_scattered_open}).
$\mathtt{color}$ and $\mathtt{note\text{-}pitch}$ are the Tier-1 sorts of
Table~\ref{tab:admissibility}; they gather like any other sort, and
their Mix is a domain reading rather than inert geometry.
Mix on $(\mathtt{num},\mathtt{num})$ is the clock trap of
Section~\ref{subsec:clock_encoding}: it is geometrically defined and
type-refused. $^{\dagger}$Mix on $(\mathtt{word},\mathtt{word})$ is
allowed as geometry and carries no linguistic $\chi$-claim
(\hyperref[subsec:axiom0_admissibility]{Axiom~0}, Tier~0).
$^{*}$Numeric $T$ and numeric offset may be defined later as
maps of the chosen numeric encoding; they are not the lexical
operators of Sections~\ref{sec:transformation_algebra}
and~\ref{sec:analogical_offset}. Offset on the Tier-1 sorts is
untested here and marked open rather than claimed.
$^{\ddagger}$Gather is sort-blind: the $\oplus$ column is \emph{yes}
for every pair, including pairs of different sorts and sorts not
listed here. A bowl records which units are present, and presence
carries no domain claim, so there is nothing for a sort mismatch to
corrupt. What gather does \emph{not} do is combine: $2\oplus
\texttt{two}=\{2,\texttt{two}\}$ is two seats and no cast, and
$1\oplus 1=\{1\}$ is one seat and not the sum $2$. Counting stays with
$+$ (Section~\ref{subsec:clock_encoding}). The final row is the
general rule; a sort seated on the sphere may always be gathered, while
its other operators wait on
\hyperref[subsec:axiom0_admissibility]{Axiom~0}. A sort inadmissible for
Tier~1, such as $\mathtt{snowflake}$ in
Table~\ref{tab:admissibility}, still falls under this general row at
Tier~0: it may be gathered, and its other operators remain open until
a stated map is given --- the same status Table~\ref{tab:admissibility}
already grants it, not a stricter one. What has genuinely not been
done for $\mathtt{snowflake}$, unlike $\mathtt{num}$, $\mathtt{word}$,
$\mathtt{color}$, or $\mathtt{chladni}$, is the construction of an
explicit chart; Tier-0 admission is asserted here, not built. A sort
with \emph{no} row in either table --- one this paper has not
considered at all --- is undecided in the stronger sense of
Table~\ref{tab:admissibility}'s closing line, and is not seated until
one is given.
$^{\S}$Mix on $(\mathtt{chladni},\mathtt{chladni})$ is a domain
reading only within one mode, where the place restriction above is
satisfied automatically; across modes it is undefined rather than
merely untested (Section~\ref{subsec:chladni_inclusion}). $T$ is
marked open because direction carries a mode label for this sort and
no domain meaning is claimed for rotating it.}
\label{tab:type_table}
\end{table}

\subsection{Worked mixed sentence}
\label{subsec:mixed_sentence}
Take the report line \emph{Wind of $12\,\mathrm{m/s}$ darkened the sun.}
The register is used as follows.

\begin{enumerate}
    \item Load $\mathrm{wind}$, sort $\mathtt{word}$.
    \item Load $12$, sort $\mathtt{num}$. Juxtapose; do not form
          $\mathrm{wind}\oplus 12$.
    \item Load $\mathrm{m/s}$ as a unit-name (lexical qualifier of the
          number).
    \item Load $\mathrm{sun}$, sort $\mathtt{word}$.
    \item Let the number cap a strength,
          $\sigma(12)\in[0,1]$, and set
          $r_{\mathrm{wind}}=\sigma(12)$ so that
          $\alpha=\pi\sigma(12)$ in Equation~\eqref{eq:transform_def}.
    \item Apply $T_{\mathrm{wind}}(\mathrm{sun})$. Output sort:
          $\mathtt{word}$.
\end{enumerate}
The scalar never became a hue. It turned a dial of an action whose
actors remained words. That is the same division of labor a printed
paper already uses: digits set magnitudes; words name what the
magnitudes belong to.

\subsection{Several lexicons, one concept sphere}
\label{subsec:shared_lexicons}
Surface tokens are not meanings. Two dictionaries may write the same
pair of concepts with unrelated strings:
\begin{equation}
    \begin{aligned}
        L_1&: \quad \texttt{xyz}\mapsto\mathrm{man},\quad
                    \texttt{qyz}\mapsto\mathrm{woman},\\
        L_2&: \quad \texttt{fgh}\mapsto\mathrm{man},\quad
                    \texttt{hgf}\mapsto\mathrm{woman}.
    \end{aligned}
    \label{eq:two_lexicons}
\end{equation}
In the usual pipeline each lexicon is tokenized, encoded, and
\emph{embedded separately}. The geometry that makes
$\mathrm{man}$ near $\mathrm{woman}$ is then learned twice, and a
further alignment step is required before $L_1$ and $L_2$ can share
an analogy. The token-to-embedding map is doing two jobs at once:
naming a string, and inventing a place for its meaning.

The unit sphere is offered as a \emph{proposed common coordinate register}, not as a demonstrated universal semantic space and not as a replacement for a lexicon table.
A token is a surface form. A state on $\mathcal{S}$ is a meaning.
Each lexicon is only a cast
\begin{equation}
    E_{L}(\texttt{token}) = x\in\mathcal{S},
    \label{eq:lexicon_cast}
\end{equation}
the same kind of map as $E_{\mathtt{num}\to\mathtt{word}}$ above.
Under that discipline both $\texttt{xyz}$ and $\texttt{fgh}$ load the
\emph{same} state $\mathrm{man}$. The offset
$\mathrm{Sun}\to\mathrm{Day}$ of Section~\ref{sec:analogical_offset}
is then computed once, on the globe, and reused by every lexicon that
points at those two points.

Embeddings therefore stay. They do not disappear.
What changes is that they are no longer private vector spaces, one
per lexicon. They are embeddings \emph{into the same unit}
$\mathcal{S}$. Once $E_{L_1}(\texttt{xyz})=E_{L_2}(\texttt{fgh})=\mathrm{man}$,
the two surface forms occupy one address, and every angle measured
from that address is the same angle:
\begin{equation}
    \angle(\mathrm{man},\mathrm{woman})
    \;=\;
    \angle\bigl(E_{L_1}(\texttt{xyz}),\,E_{L_1}(\texttt{qyz})\bigr)
    \;=\;
    \angle\bigl(E_{L_2}(\texttt{fgh}),\,E_{L_2}(\texttt{hgf})\bigr).
    \label{eq:uniform_angles}
\end{equation}
Under that proposal the embeddings would be uniform: same globe, same points for the same
concepts, same angles for the same relations. A raw string still
needs the lexicon map $E_L$. It would no longer need a new geometry.
The algebra does not mint a universal $\mathrm{man}$ point by itself. That is a dictionary programme, not an axiom.

\subsection{Deferred programme: standardized concept dictionaries}
\label{subsec:standard_dictionaries}
The present section only \emph{names} the shared globe. It does not
build the dictionaries, and it does not treat them as part of the
algebra closed in Axioms~0--7. The later refinement, if the idea is
pursued, is a pair of catalogues rather than a new operator, and is
filed under Future Work (Section~\ref{sec:conclusion}):
\begin{itemize}
    \item one \emph{concept dictionary} --- a stipulated inventory of
          unit states on $\mathcal{S}$ (or a higher-dimensional
          analogue), so that $\mathrm{man}$ and $\mathrm{woman}$ have
          standard addresses and therefore standard angles;
    \item one \emph{language lexicon} $E_L$ per written or spoken
          system, sending that language's tokens onto those addresses;
    \item learnable paths --- the species of $T$ and the dials inside
          a given $T$ (angle map, phase map, paint-versus-poem gate)
          may be fitted from examples. Learning would adjust the path,
          not the rule that a bowl is collected before it is read.
          See also Section~\ref{sec:conclusion}.
\end{itemize}
Standardization of the first catalogue is what makes embeddings
uniform across English, a cipher, or a second natural language.
The second catalogue is still a tokenizer-plus-encoding: it is the
phone book, not the city. Tokens and encodings have already earned
their keep, so they stay. They only have to work \emph{lighter}:
look up a standard address instead of inventing one. They disappear as
\emph{inventors of geometry}. They do not disappear as \emph{names}.
A world in which every language pointed at the same concept
dictionary would still print letters on a page; It would no longer need each 
corpus to learn a private $\mathrm{man}$--$\mathrm{woman}$ angle. 
That programme is recorded here and left open. It is not an
axiom, and it is not implied by the existence of $\mathcal{S}$ alone.

\subsection{Scope of the register model}
\label{subsec:register_scope_final}
This section claims only that one sphere can \emph{host} both alphabets
as a workspace, and that several lexicons can point at the same
workspace through different casts. It does not claim that the two-dimensional surface
$S^2$ is a realistic embedding space for a full natural-language
lexicon: a working word embedding typically needs tens or hundreds of
dimensions. The construction here is a typed register and a toy
dictionary on one globe. Scaling the lexical encoding is a later
choice --- a higher-dimensional sphere, a spherical-harmonic field, or
an external lookup table that writes states into the register. The
lookup tables $E_L$ may still be learned; the claim is only that they
should write into one geometry rather than each minting a private one.
The unit ball of radius $1$ in $\mathbb{R}^3$ is where the idea is
taught. The unit sphere in the embedding dimension is where close
semantics can live. The algebra --- collapse, path, result --- is
meant to travel with the dimension. It is not meant to replace $768$
numbers with two angles and hope the library still sorts.

\section{A Sketch of Conventions}
\label{sec:conventions_sketch}
These notes are a seed for later grammars, not a code to enforce.
They show one way a typed register \emph{might} be spoken. A later
grammar --- learned, cultural, or formal --- may replace every line.
This paper only needs the ground level enough that such a grammar has
somewhere to stand.

\paragraph{Case does not move the address.}
The tokens $\texttt{man}$, $\texttt{Man}$, and $\texttt{MAN}$ load the
same direction $(\theta,\phi)$. Orthography is a doorbell, not a new
city block.

\paragraph{Language does not move the address.}
English $\texttt{man}$ and a Polish token for the same concept load the
same point. Only the lexicon map $E_L$ changes. Language is a
phone-book tag $L$, not a second angle.

\paragraph{Color is reserved for mood, not for language.}
Phase $\chi$ is already used as affect in the poem example
(Section~\ref{subsec:transform_example}). If language or case also
occupied $\chi$, that example would break. Emphasis may tint $\chi$
or raise $r$; it should not relocate $\mathrm{Sun}$ to $\mathrm{Moon}$.

\paragraph{A bag is not a sentence.}
$\{\mathrm{sun},\mathrm{wind}\}$ is $\oplus$.
``Wind darkens the sun'' is $T$. Ingredients and order are different
verbs. How often a unit is used is not a second shell in the bag:
the path may visit $\texttt{twinkle}$ twice; $\mathtt{num}$ may
count $1+1+1+1=4$. $\oplus$ will not.

\paragraph{A numeral is not a word until it is cast.}
The digit $2$ stays $\mathtt{num}$. The word $\texttt{two}$ is
$\mathtt{word}$. They may share a bowl --- gather is sort-blind ---
but $2\oplus\texttt{two}=\{2,\texttt{two}\}$ is two seats, not a cast
and not a sum. Becoming a word requires the explicit
$E_{\mathtt{num}\to\mathtt{word}}$ of
Section~\ref{sec:encodings_sorts}, which no gather performs.

A toy row makes the first three conventions visible. Place
$\mathrm{man}$ at $(\theta,\phi)=(10^\circ,40^\circ)$:

\begin{center}
\begin{tabular}{llcc}
Token & $L$ & $(\theta,\phi)$ & $\chi$ (mood only) \\
\hline
\texttt{man} & English & $(10^\circ,40^\circ)$ & $0^\circ$ (neutral) \\
\texttt{Man} & English & $(10^\circ,40^\circ)$ & $0^\circ$ \\
\texttt{MAN} & English & $(10^\circ,40^\circ)$ & $20^\circ$ (emphasis, optional) \\
\texttt{man} & Polish  & $(10^\circ,40^\circ)$ & $0^\circ$ \\
\end{tabular}
\end{center}

Same block. Different doorbells. Nothing here is binding.

\section{Operational Algebra \& Numerical Demonstration}

\subsection{Mix of two shells}
Collecting $r=0.6$ at $0^\circ$ with $r=0.8$ at $120^\circ$ (same
place) yields two shells, not an angle. Applying Mix ($\boxplus$) and
Equation~\eqref{eq:interference_formal} yields the 4-tuple
$(r',\theta',\phi',\chi')\approx(0.721,\,\theta,\,\phi,\,73.9^\circ)$.
The older figure $66.6^\circ$ is withdrawn: Axiom~5 as a length-only
formula never produced it.

\subsection{Rotational transposition, and what division is not}
\label{subsec:division_open}
Rotation is the rule of \S\ref{subsec:axiom_rotation}: $\mathbf{R}\in\mathrm{SO}(3)$
turns $\hat n$. Inverse rotation is well-defined. Division as
``deconvolution and signal unmixing'' is not defined in this paper
and is left open --- shown here to be ill-posed for this algebra
specifically, not merely unbuilt.

Given only a Mixed result, the bowl that produced it is not merely
hard to recover; it is not determined by the result at all. The
introduction's own red-and-green example, $0^\circ$ and $120^\circ$
at equal strength reaching yellow at $60^\circ$, is one point in an
infinite family: two components at $10^\circ$ and $100^\circ$,
correctly weighted, reach the identical phasor sum, to machine
precision --- verified directly, not merely argued, and neither pair
is preferred by the result over the other. Assuming a specific
composition and solving for it does not rescue this. A wrong assumed
composition still reconstructs the observed result exactly: solving
for $10^\circ$-and-$100^\circ$ strengths against the true
red-and-green result returns a residual of $10^{-17}$, indistinguishable
by any check internal to the algebra from the correct pair's own
residual of zero. An algebra that states only forced, verified
identities has no honest place for an operation whose wrong answer
looks exactly like its right one.

This is deliberately left to inference (companion paper,
\emph{Pattern Arithmetic: Inference from Bowls and Paths}), where a
comparison against a witnessed reference is held to a stated
confidence and may abstain --- a discipline this algebra's own Mix
does not have available to it, and should not manufacture.

\subsection{Spectral storage (optional, later)}
Spherical-harmonic arrays of size $O(L^2)$ may later store a field of
many atoms. They are not the working atom of
\hyperref[subsec:axiom1]{Axiom~1}, and $O(N_{\mathrm{vox}}^3)\to O(L^2)$ is not a theorem of
Pattern Arithmetic as it stands.

\section{Higher-order systems as bowls and paths}
\label{sec:higher_order}
These examples illustrate the $\oplus$-versus-$T$ split with borrowed
vocabulary. They are costumes
(Section~\ref{subsec:claim_scope}), not models of weather, ecology,
or disease.
Higher-order truths are bowls of basic units read by a path.
A storm, a forest, or a cancellation is such a reading.
It is an example of the spine, not a model of the sky, the woods,
or the clinic. The path defines the higher-order system: the bowl
only lists who is present.

A vertical pinball table makes the same split visible. Each bumper
is one seated unit --- a direction and a size \(r\). The set of
bumpers on the board is the bowl (wind, temperature, humidity,
pressure). The pockets at the bottom are dictionary pins already
named \(\mathrm{storm}\) and \(\mathrm{calm}\). Dropping a ball from
a slot is a path \(T_{\mathrm{weather}}\): the ball meets some
bumpers, in some order, and each kick depends on that bumper's
angle and strength. The landing pocket is a comparison after the
walk. Adding the bumpers together does not \emph{make} a pocket.
A different slot, or a different order of kicks, is a different
weather. This is one rough picture of \emph{a} transform path, not
the mechanics of every $T$, and not a second axiom. A Mix reading ($\boxplus$), and a poem path $T_{\mathrm{poem}}$, need not bounce.

\subsection{Atmospheric storm (costume): one pipeline with numbers}
\label{subsec:storm_pipeline}
This is a costume for the spine, not a model of the sky.
These walks are not weather equations, and they are not forest
equations. They are still computations: ordered readings of many
co-present units that can land near a higher-order name already
seated in a dictionary.
Four incoming
states, one bowl, two paths, two endpoints. Every figure below is
Rodrigues' rotation of $\hat n$ by $\alpha=\pi r$ about the actor's
axis, with $\chi$ added. A reader can recompute.

\paragraph{1. Input states.}
\begin{align*}
\mathrm{Wind}&:
  \; r=1,\;
  \hat n=(1,0,0),\;
  \chi=0^\circ,\;
  \tau=\mathtt{word},
  \quad \alpha=180^\circ,\\
\mathrm{Heat}&:
  \; r=\tfrac12,\;
  \hat n=(0,1,0),\;
  \chi=90^\circ,\;
  \tau=\mathtt{word},
  \quad \alpha=90^\circ,\\
\mathrm{Moisture}&:
  \; r=\tfrac23,\;
  \hat n=(0,0,1),\;
  \chi=45^\circ,\;
  \tau=\mathtt{word},
  \quad \alpha=120^\circ,\\
\mathrm{Eye}&:
  \; r=1,\;
  \hat n=\bigl(\tfrac{\sqrt{2}}{2},\,0,\,\tfrac{\sqrt{2}}{2}\bigr),\;
  \chi=0^\circ,\;
  \tau=\mathtt{word}.
\end{align*}
The eye sits at $(\theta,\phi)=(45^\circ,0^\circ)$.

\paragraph{2. Collapse into the bowl.}
\begin{equation}
    R
    \;=\;
    \mathrm{Wind}\oplus\mathrm{Heat}\oplus\mathrm{Moisture}\oplus\mathrm{Eye}
    \;=\;
    \{\mathrm{Wind},\,\mathrm{Heat},\,\mathrm{Moisture},\,\mathrm{Eye}\}.
    \label{eq:storm_bowl}
\end{equation}
Four distinct names, four seats. A second $\mathrm{Wind}$ would not
open a fifth. Call this bowl $s(\mathrm{weather})$ if a name helps:
it is still four pins, not a fifth vector. Mix is a different
reading and is not the storm.
In the dictionary, $\mathrm{storm}$ and $\mathrm{calm}$ are
\emph{other} seated units --- separate doorbells, not sums of
wind, temperature, humidity, and pressure. Mere $\oplus$ cannot
reach those pins. $T_{\mathrm{weather}}$ is the path family that
reads $s(\mathrm{weather})$; $T_{\mathrm{deepen}}$ and
$T_{\mathrm{clear}}$ below are two such walks. An endpoint may
fall \emph{near} the storm pin or the calm pin. Nearness is
comparison after the walk, not addition of the parts.

\paragraph{3. Two trajectory paths.}
The register occupant is the eye. The other three shells act, in order:
\begin{equation}
    T_{\mathrm{deepen}}
    \;=\;
    T_{\mathrm{Wind}}\circ T_{\mathrm{Heat}}\circ T_{\mathrm{Moisture}},
    \qquad
    T_{\mathrm{clear}}
    \;=\;
    T_{\mathrm{Moisture}}\circ T_{\mathrm{Heat}}\circ T_{\mathrm{Wind}}.
\end{equation}

\emph{Deepen} --- moisture, then heat, then wind:
\begin{align*}
T_{\mathrm{Moisture}}(\mathrm{Eye})
&= \bigl(1,\;(-\tfrac{\sqrt{2}}{4},\,\tfrac{\sqrt{6}}{4},\,\tfrac{\sqrt{2}}{2}),\;45^\circ\bigr),\\
T_{\mathrm{Heat}}(\cdots)
&= \bigl(1,\;(\tfrac{\sqrt{2}}{2},\,\tfrac{\sqrt{6}}{4},\,\tfrac{\sqrt{2}}{4}),\;135^\circ\bigr),\\
T_{\mathrm{Wind}}(\cdots)
&= \bigl(1,\;(\tfrac{\sqrt{2}}{2},\,-\tfrac{\sqrt{6}}{4},\,-\tfrac{\sqrt{2}}{4}),\;135^\circ\bigr).
\end{align*}

\emph{Clear} --- wind, then heat, then moisture:
\begin{align*}
T_{\mathrm{Wind}}(\mathrm{Eye})
&= \bigl(1,\;(\tfrac{\sqrt{2}}{2},\,0,\,-\tfrac{\sqrt{2}}{2}),\;0^\circ\bigr),\\
T_{\mathrm{Heat}}(\cdots)
&= \bigl(1,\;(-\tfrac{\sqrt{2}}{2},\,0,\,-\tfrac{\sqrt{2}}{2}),\;90^\circ\bigr),\\
T_{\mathrm{Moisture}}(\cdots)
&= \bigl(1,\;(\tfrac{\sqrt{2}}{4},\,-\tfrac{\sqrt{6}}{4},\,-\tfrac{\sqrt{2}}{2}),\;135^\circ\bigr).
\end{align*}

\paragraph{4. Endpoints.}
$T$ leaves the occupant's strength at $r=1$. Phase addition commutes, so
both paths end at $\chi=135^\circ$. The directions do not agree:
\begin{align*}
T_{\mathrm{deepen}}(\mathrm{Eye}):
\quad
\hat n
&= \bigl(\tfrac{\sqrt{2}}{2},\,-\tfrac{\sqrt{6}}{4},\,-\tfrac{\sqrt{2}}{4}\bigr)
\approx (0.7071,\,-0.6124,\,-0.3536),\\
(\theta,\phi)&\approx(110.7^\circ,\,319.1^\circ);\\[0.4em]
T_{\mathrm{clear}}(\mathrm{Eye}):
\quad
\hat n
&= \bigl(\tfrac{\sqrt{2}}{4},\,-\tfrac{\sqrt{6}}{4},\,-\tfrac{\sqrt{2}}{2}\bigr)
\approx (0.3536,\,-0.6124,\,-0.7071),\\
(\theta,\phi)&=(135^\circ,\,300^\circ).
\end{align*}
Same bowl, two weathers. The storm is the walk. It is not a fifth pin,
and it is not a claim about vorticity or shear.

\subsection{Table and chair (costume)}
\label{subsec:table_chair}
A table is one seated unit; a chair is another. $\{\mathrm{table},\,\mathrm{chair}\}$
is their bowl under $\oplus$ --- two names present, order-blind, no
claim about arrangement. A room is not read from that bowl. It is
read from a path: $T_{\mathrm{chair}}(\mathrm{table})$, the chair
acting on the table, is a placement --- pulled up to the table's edge.
$T_{\mathrm{table}}(\mathrm{chair})$, the same two shells with the
actor swapped, is a different placement --- the table built around
a fixed chair. $\oplus$ cannot tell these apart; it only knows both
are present. $T$ can, because $T$ is not symmetric in its two
arguments the way $\oplus$ is. As with the storm's pockets, a named
room --- ``dining nook,'' ``reading corner,'' ``cluttered storage''
--- is a separate dictionary pin, not a sum of the furniture. An
endpoint may land near one pin or another; nearness is comparison
after the walk, exactly as in Section~\ref{subsec:storm_pipeline}.
This is a costume for the same split table and chair make familiar
without any physics: two objects merely being present is one
reality (grocery counting: one bag of grain plus one bag of grain is
two bags --- schoolbook $+$, not $\oplus$); two objects placed
in a specific relation is a different reality, and no amount of
$\oplus$ on the parts reaches it. It is not a model of furniture,
interior design, or how rooms are actually named.

The furniture case is also where the motivating hypothesis behind
this whole paper is easiest to say plainly: a path $T_{\mathrm{chair}}$
that keeps ``chair acts on table'' as one typed, invertible step is a
shorter thing to learn than the same fact buried inside a flat
embedding, where the model must first discover, from undifferentiated
vectors, that arrangement is a different kind of fact than mere
co-occurrence. Keeping presence ($\oplus$) and path ($T$) as separate
primitives is a bet that higher-order truths become easier to learn
because the distinction is built in rather than re-derived from data.
That bet is stated here as a hypothesis, not a result: no example in
this paper demonstrates it, and one is deliberately not attempted.
Section~\ref{subsec:future_work} leaves it open.

\subsection{Forest (costume)}
Flora, fauna, soil, moisture sit as units. $\oplus$ gathers them into
one bowl --- four names present, order-blind
(\hyperref[subsec:axiom_idempotence]{Axiom~4}).
There is no $A+B=AB$. A path may read competition or succession.
The forest-as-system is that reading.

\subsection{Pathogen cancellation (costume)}
A mark and an equal-strength anti-phase mark gather as two seats in
one bowl (\hyperref[subsec:axiom_idempotence]{Axiom~4}); $\oplus$ does not
cancel them. That bowl may then be \emph{read} by Mix into black
(\hyperref[subsec:axiom_interference]{Axiom~5}). That is geometry. It is not
medicine.

\section{The Enclosure: Seats, Rest, and Camps}
\label{sec:enclosure}

Every operator before \hyperref[subsec:axiom_assembly]{Axiom~6} works
entirely on the bounded stage: units go in, units or states come
out, and nothing asks where anything sits. Assembly asks. This
section supplies the room it asks about, and then follows what
happens when units are left alone in it.

The enclosure is a costume borrowed briefly and then dropped. Many
small things, left alone, falling into a few quiet piles is a
familiar picture from astronomy; the units here are not stars, carry
no mass, and move only by the rules below. After this paragraph the
working words are \emph{seat}, \emph{rest}, and \emph{camp}.

\subsection{Seats}
\label{subsec:enclosure_seats}
A unit, or a bowl, may carry a \emph{seat}: a position vector
$\vec P$ measured from one shared origin, in ordinary unbounded
Euclidean space. The seat is external to the 4-tuple of
\hyperref[subsec:axiom1]{Axiom~1}, in the same way step metadata sits on the
path rather than inside the units it walks. The unit is unchanged by
having one; $r$ is still bounded, $\hat n$ is still a semantic
address and not a physical one. For two seated units,
\begin{equation}
    d_{AB} = \lvert \vec P_B - \vec P_A\rvert,
    \qquad
    \hat r_{AB} = \frac{\vec P_B - \vec P_A}{d_{AB}},
    \label{eq:enclosure_distance}
\end{equation}
which completes Equation~\eqref{eq:assembly_energy}. Where seats
exist, $\hat r_{AB}$ is the separation axis; no axis derived from
orientation is required or used.

\subsection{Rest}
\label{subsec:enclosure_rest}
Seats may move; meaning does not. $\hat n$ and $\chi$ do not crawl,
so the lexicon stays put and no unit drifts into being a different
unit. What moves is $\vec P$, and it moves by descending the energy
of Equation~\eqref{eq:assembly_energy}:
\begin{equation}
    \vec P_i \;\leftarrow\; \vec P_i + \delta\,F_i,
    \qquad
    F_i = -\nabla_{\vec P_i}\!\!\sum_{j\ne i} U_{ij}.
    \label{eq:enclosure_motion}
\end{equation}
There is no mass term and no momentum carried between steps. This
is deliberate rather than an omission. A unit carries no inertia, so
it cannot overshoot: an isolated pair closes distance monotonically
and never exceeds its own starting separation. Weighting each step
by a mass-like quantity was tested and produces oscillation rather
than rest, precisely where the weighting is most extreme --- a
ninefold difference in strength gives a $729$-fold weighting under a
cubed reading, and a demonstrably non-monotonic trajectory against
the clean decline of Equation~\eqref{eq:enclosure_motion}
(the printed pair of \S\ref{subsec:enclosure_replay}).

A direct consequence: since $F_A=-F_B$ exactly and nothing divides
the resulting step, two units of any polarity move by identical
magnitude toward each other. Polarity labels a bond; it does not
sort the movers. A second consequence: since every force comes in an
equal and opposite pair, the plain, unweighted mean of the seats is unchanged by every step.

\paragraph{Rest is a freeze, not a stop.}
Taken literally, rest arrives only once, and only at the end. The
pull between two units of nonzero strength is proportional to
$d\,(d^2+\varepsilon^2)^{-5/2}$, which is positive at every
separation $d>0$: no two such units are ever exactly still while
apart. Given unlimited time, every group eventually joins every
other and all seats of nonzero strength meet at one point. What the enclosure offers
instead is a stated freeze. A configuration is \emph{at rest at
resolution $\eta$} when
\begin{equation}
    \max_i \lvert F_i \rvert \;<\; \eta,
    \label{eq:enclosure_freeze}
\end{equation}
and a configuration that has reached this threshold may be frozen and
taken as the considered one --- the accepted solution at that
resolution. Its shape, the walk that produced it, and its camps are
then what is read. The threshold $\eta$ is a measurement convention,
named alongside the step $\delta$ and the census cut below, and not a
truth of the algebra: a finer $\eta$ runs longer and may see groups
merge that a coarser one froze apart. It is stated on force rather
than on the size of a step for a reason. A step is $\delta F_i$, so a
threshold on steps would let the choice of $\delta$ decide which
configurations count as rest; a threshold on force does not.

The step $\delta$ in Equation~\eqref{eq:enclosure_motion} is a
numerical convention and not a truth of the algebra. It sets how
finely the descent is sampled, not where the freeze falls. For one
isolated group this is exact: the plain mean of its seats is
conserved, so it comes to rest on the plain mean of its starting
seats whatever $\delta$ is used, agreeing across
$\delta = 0.02$, $0.01$, $0.005$, and $0.002$ to within
$10^{-15}$. Where several groups pull on one another weakly, the
same four steps reach the freeze at the same time and with the same
camps, and their bodies agree to within $2\times10^{-7}$ --- the
ordinary discretization error of sampling the slow drift between
groups (\S\ref{subsec:enclosure_replay}). A step chosen too coarsely
relative to $\varepsilon_{AB}$ will misbehave near close approach,
in the ordinary way of any fixed-step scheme; that is a limit on the
arithmetic, not on the law.

Inside a group, the merge does not stop short of a point. Two
members close together pull in proportion to their separation, so a
frozen group's members sit on one seat to within a distance that
shrinks with $\eta$ --- under $2\times10^{-5}$ at $\eta=10^{-5}$ in
the configuration of \S\ref{subsec:enclosure_replay}. A frozen camp
has a location and a membership. It has no internal shape.

\subsection{Camps and bodies}
\label{subsec:enclosure_camps}
Units left alone under Equation~\eqref{eq:enclosure_motion} do not
always fall into one pile. The energy dies as $d^{-3}$, so distant
units are nearly irrelevant to one another and local groups settle
before they feel the far ones. Several separated groups can
therefore reach the freeze of Equation~\eqref{eq:enclosure_freeze}
together, each one a \emph{camp} --- a camp at resolution $\eta$,
since without a freeze they would go on to merge. The number of
camps is not chosen in advance; it is whatever the law, the starting
seats, and $\eta$ produce.

For each camp $G_k$ that has frozen, its \emph{body} is
\begin{equation}
    \vec P_{c,k}
    = \frac{\sum_{i\in G_k} r_i\,\vec P_i}{\sum_{i\in G_k} r_i},
    \label{eq:enclosure_body}
\end{equation}
the strength-weighted centre of its members. Nothing computes this
during motion and no unit seeks it; each unit feels only the sum of
pairwise pulls from its neighbours. The body is measured afterwards,
by an observer, and is a description rather than a destination.

Where units carry inner structure, the weights $r_i$ in
Equation~\eqref{eq:enclosure_body} may be replaced by
$\lvert V_i\rvert$, the magnitude of a unit's own resultant. A unit
whose inner arrows cancel still sits in the camp but barely moves its
body --- the body then reports surviving configuration rather than
bare presence. Both readings are legitimate; they answer different
questions and should not be conflated. The difference is a property
of a camp read \emph{before} its freeze. Once frozen, a camp's
members sit on one seat to the resolution of $\eta$, so every
weighting returns the same body to that resolution, and the body of
an isolated camp is simply the plain mean of its starting seats.

A system of several camps is described by the \emph{set} of bodies
$\{\vec P_{c,1},\vec P_{c,2},\ldots\}$, never by one grand average
over all of them. Two camps ten units apart have no member at their
midpoint, and reporting one body there would assert a camp that does
not exist. This is the same honesty the bowl already practises when
Mix returns black: where the evidence is two truths, do not mint a
third that nothing occupies.

Counting camps requires a stated separation cut, or a chain of short
links between members. That cut is a measurement convention and not
a truth of the algebra: the camps are produced by the law, the
census of them is not.

\paragraph{A camp is a checkpoint, not a destination.}
Every camp is read at some freeze $\eta$, after some elapsed time; a
longer wait, a finer $\eta$, or a different $\delta$ freezes a
different, though usually nearby, configuration. Call each such
freeze a \emph{checkpoint}. Two checkpoints of the same starting
seats --- upstream or downstream in time, at a different $\eta$, or
under a different $\delta$ --- need not agree exactly, and
\S\ref{subsec:enclosure_replay} already prints the drifts: $2\times
10^{-7}$ between step sizes, $2.4\times10^{-8}$ between weightings,
$3\times10^{-5}$ across a weak inter-group pull. None of these is a
mistake; each is what the same law gives at a different checkpoint.

Treating two checkpoints as \emph{the same} freeze therefore
needs one further named convention, an \emph{allowance}
$\epsilon_{\text{allow}}$: checkpoints whose bodies differ by less
than $\epsilon_{\text{allow}}$ are read as one, and checkpoints that
differ by more are read as genuinely apart. This is not a new law.
It sits beside $\eta$, $\delta$, and the census cut as a fourth
measurement convention, stated once and applied consistently, not a
truth of Equation~\eqref{eq:assembly_energy}. Every drift printed
above falls well inside a stated allowance of, say, $10^{-4}$; the
storm's Deepen and Clear endpoints of \S\ref{subsec:storm_pipeline},
$0.500$ apart in direction ($28.9^\circ$), do not fall inside any
allowance a reader would state honestly. An allowance too generous
erases real structure; one too strict reports a new camp
every time the integrator takes a different step. Naming it is the
price of being able to say two checkpoints are the same thing at
all.

\subsection{A replayable configuration}
\label{subsec:enclosure_replay}
The claims of this section are printed against one configuration, so
that any reader can replay them. How a configuration is generated ---
by hand, or by any library under any seed --- is not part of the
claim; the printed seats are. Six units in two groups forty apart:

\begin{center}
\begin{tabular}{cccc}
Unit & $r$ & $\hat n$ & starting $\vec P$ \\
\hline
$1$ & $0.5$ & $(1,0,0)$ & $(0,0,0)$ \\
$2$ & $0.6$ & $(0,1,0)$ & $(1,0,0)$ \\
$3$ & $0.7$ & $(0,0,1)$ & $(0,1,0)$ \\
$4$ & $0.8$ & $(1,0,0)$ & $(40,0,0)$ \\
$5$ & $0.4$ & $(\tfrac{1}{\sqrt2},\tfrac{1}{\sqrt2},0)$ & $(41,0,0)$ \\
$6$ & $0.9$ & $(0,0,1)$ & $(40,1,0)$ \\
\end{tabular}
\end{center}

Under Equation~\eqref{eq:enclosure_motion} with the freeze
$\eta=10^{-5}$ of Equation~\eqref{eq:enclosure_freeze}, the runs at
$\delta=0.02$, $0.01$, $0.005$, and $0.002$ take $628$, $1261$,
$2526$, and $6320$ steps --- the same elapsed time, $12.6$, in each
--- and every one freezes into the same two camps, units $\{1,2,3\}$
and $\{4,5,6\}$. The bodies land at
$(0.33336,\,0.33333,\,0)$ and $(40.33330,\,0.33333,\,0)$: the plain
means of each group's starting seats, $(\tfrac13,\tfrac13,0)$ and
$(\tfrac{121}{3},\tfrac13,0)$, each drawn about $3\times10^{-5}$
toward the other by the weak pull across the gap. The four runs agree
on both bodies to within $2\times10^{-7}$; the plain mean of all six
seats is unchanged throughout; and the $r$-weighted and plain bodies
of the first camp differ by $2.4\times10^{-8}$. Run alone, the first
group comes to rest on $(\tfrac13,\tfrac13,0)$ at all four steps,
agreeing to within $10^{-15}$. Neither camp is still in the
literal sense: the two remain forty apart only because their mutual
pull fell below $\eta$ before they could close the distance.

The mass-weighting comparison of \S\ref{subsec:enclosure_rest} is
replayed on a printed pair: $r=0.9$ at $(0,0,0)$ and $r=0.1$ at
$(3,0,0)$, both with $\hat n=(1,0,0)$, $\delta=0.02$, $3000$ steps.
Under Equation~\eqref{eq:enclosure_motion} the separation never
increases, falling from $3$ to $1.80$. Dividing each step by $r^3$ ---
a $729$-fold weighting between the two --- makes the separation
increase on $804$ of the $3000$ steps.

A printed configuration makes a result replayable. It does not make
the result right: a replay reproduces what the law did from those
seats, and whether that answer is the one a reader wanted is a
separate question.

\subsection{What the enclosure does not supply}
\label{subsec:enclosure_open}
Assembly never repels, so nothing in the enclosure pushes members
apart; separated camps persist because pull dies with distance, not
because anything resists, and only until a freeze is declared. There
is no final plurality: left without a freeze, every camp eventually
joins every other. There is correspondingly no core, no halo,
and no sorting of members by strength --- structures familiar from
the borrowed picture that this law does not reproduce, and should not
be graded against. A camp does not acquire a name by coming to rest.
Naming is comparison against a seated pin, which is a different
question asked afterwards and elsewhere.


\section{Discussion \& Related Work}
\label{sec:related}

Each paragraph below states what is reused and what is added. None of
these citations is a claim of identity. The list is a
boundary-drawing exercise against formalisms that share a sphere, an
angle, or a rotation
(Section~\ref{subsec:claim_scope}). It is not a unification of those
fields, and it is not a proposal to replace quantum mechanics,
embeddings, or transformers. Empty headings are omitted.

\paragraph{Bloch sphere.}
A pure qubit is a point on $S^2$; a mixed qubit lives in the ball.
Pattern Arithmetic reuses that unit globe. It differs in three
respects. Strength $r\le 1$ is an activation of one arrow, not a purity
parameter of a density matrix. Phase $\chi$ is an internal colour of
that arrow, not a global $U(1)$ that can be quotiented away. There is
no measurement postulate: Mix may mint a third unit or black; $\oplus$
only collects. They do not collapse a superposition onto an eigenstate. Related work
may say: we reuse a unit globe; we do not borrow quantum authority for
$\oplus$.

\paragraph{Word2vec and cosine embeddings.}
Lexical similarity as an angle is reused. A bag of words as a
commutative gather is reused in spirit. What is added is a bounded unit
($r\le 1$), an explicit collection operator that is idempotent
($A\oplus A=A$), and an order-dependent path $T$ that cosine addition
does not supply. Offset is a manifold cousin of the familiar
$\mathrm{king}-\mathrm{man}+\mathrm{woman}$ arithmetic, not a proof
that the toy globe is a full lexicon.

\paragraph{Transformers.}
A transformer occupies the same \emph{slot} as $T$: it is a path
through seated tokens. It is not the same device. The usual stack is
\[
\text{sentence}\to\text{token}\to\text{encode}\to\text{embed}\to\text{transformer}.
\]
The stack of this paper is
\[
\text{sentence}\to\text{token}\to E_L\to\text{sphere}\to
\oplus\text{ (bowl) }+\; T\text{ (walk)}.
\]
Tokens and encodings stay. The embedding is a cast into one sphere
(of whatever dimension the lexicon needs). A higher-order unit (a
line, a report, a lexicon) may treat each information atom as the
already-furnished stand-in for a token: that is a closure over many
units, not a change to the atom. $T_{\mathrm{poem}}$,
$T_{\mathrm{prose}}$, and $T_{\mathrm{equation}}$ are species of
walk, not extra nets. A transformer may later fill those walks as a
learned path. This paper does not claim to replace attention; it
splits presence from order so that a learned path need not also
invent the city.

\paragraph{RotatE.}
Knowledge-graph embeddings that treat a relation as a rotation in a
complex plane are the nearest neighbour of $T$. Pattern Arithmetic
extends that idea from $S^1$ (or a product of circles) to
$\mathrm{SO}(3)$ acting on $S^2$, and pairs the rotation with an
additive phase $\chi$. It also keeps a separate collapse operator,
which RotatE does not.

\paragraph{von Mises--Fisher.}
Directional statistics on the sphere reuse the same stage.
Those models estimate a density. Pattern Arithmetic does not: the atom
is a point, $r$ is not a concentration $\kappa$, and $\oplus$, $\boxplus$, $T$, and
offset are algebraic operations on units, not maximum-likelihood fits.

\paragraph{Riemann sphere.}
Stereographic compactification of $\mathbb{C}\cup\{\infty\}$ onto $S^2$
is the natural home of one numeric encoding (unbounded reals toward the
poles). The algebra of this paper is not the M\"obius group. A clock,
or a stereo map, is a chart; $+$ in $\mathtt{num}$ mode is decoded
arithmetic, not a fractional linear map.

\paragraph{$\mathrm{SO}(3)$.}
In three dimensions, the transformation $T$ is a named action of the
rotation group, plus a phase shift; beyond three, each $T$ is a simple
rotation --- one plane turned, the rest held --- and not a general
element of $\mathrm{SO}(N_{\mathrm{dim}})$
(\S\ref{subsec:transform_higher_dim}). That is the whole geometric engine of a path. What
$\mathrm{SO}(3)$ does not supply is the bowl, the cap $r\le 1$, the
sort $\tau$, or Mix. Non-commutativity of poems is inherited;
idempotence of $\oplus$ is not.

\paragraph{Spherical harmonics.}
A later field representation of many arrows could be stored as
$Y_\ell^m$ coefficients, replacing an $O(N_{\mathrm{vox}}^3)$ voxel cube by an
$O(L^2)$ spectral array. That is a compression option for a field, not
a definition of the atom, and it is not claimed as a result here.

\section{Conclusion \& Future Directions}
\label{sec:conclusion}

The claim of this paper is narrow
(Section~\ref{subsec:claim_scope}). Counting is not the only
arithmetic for combining units of information. A combination rule
should match the unit. The sphere and the operators
$\oplus$, $\boxplus$, $T$, offset, and $\bowtie$ are one proposed
encoding of that claim. They are not a general theory of language,
physics, biology, or computation. The costumes, the enclosure, and
related-work paragraphs are illustrations and boundaries, not domain
models.

What was written is two pipelines on the same unit. Collapse many
sorted units into a bowl of typed atoms $x=(r,\hat n,\chi,\tau)$ on a
unit sphere, then derive results by ordered transformation of a
register occupant. Separately, assemble units without fusing them
and let external seats descend until a stated threshold freezes
them; the frozen camps have bodies, not names. Naming a camp is
left open here, on purpose: a second companion paper,
\emph{Pattern Arithmetic: Higher Truths from Assembly}, now states
it as the question still to be answered, not a comparison already
made.
Embeddings, on this view, are
not abolished. They are embeddings into one proposed common coordinate
register, so that a
concept such as $\mathrm{man}$ \emph{may} keep the same angles under every
lexicon that points at it. That is a design, not a demonstrated
universal semantic space.

A few spoken conventions are sketched in
Section~\ref{sec:conventions_sketch} as seed, not statute.

\subsection{Future work}
\label{subsec:future_work}
What remains to be built, and is only filed here, is a standardized
concept dictionary together with per-language lexica $E_L$
(Section~\ref{subsec:standard_dictionaries}). Until those catalogues
exist, tokenization and encoding remain the handles by which a
paragraph finds its points. They may continue; they will only have to
work lighter. The open bet is that they need not also
remain the architects of a new space of meanings.
The algebra closed in this paper does not produce a universal
$\mathrm{man}$ point by itself. Standardization is sociological and
empirical; it is not an axiom.

A further deferred question is whether a path $T$ should stay
hand-specified or carry learnable weights: the angle law
$\alpha(r)$, a phase map, a gate between Mix ($\boxplus$) and
$T_{\mathrm{poem}}$, and the lexicon casts $E_L$. The honest
constraint is structural. If undo-the-poem matters, learning should
stay inside $\mathrm{SO}(3)$, not collapse $T$ into an arbitrary
matrix. Weights may change how the shells are thrown. They should
not let the shells throw themselves.

Left fully open is the hypothesis raised in
Section~\ref{subsec:table_chair}: that separating presence
($\oplus$) from path ($T$) as distinct primitives should make
higher-order truths cheaper to learn than recovering the same
distinction from undifferentiated vectors. Testing this would need a
controlled comparison --- a compositional task, a flat-embedding
baseline, and a $T$-path model, measured on data efficiency rather
than final accuracy --- and no such comparison is attempted here.
The table-and-chair and storm costumes motivate the hypothesis; they
do not test it.

Ambient $\mathbb{R}^3$ makes the expressions
$\mathcal{S}+(a,b,c)$ and $\mathcal{S}\cdot(a,b,c)$ look legal.
They are legal as \emph{vector} arithmetic. They are not automatically
legal as \emph{unit} arithmetic. A state with $r=0.98$ added to
another with $r=0.98$ along the same direction is $r=1.96$, which
has left the ball. A scalar $2\cdot(0.98)$ does the same. The unit
cap $r\in[0,1]$ then forces a choice, to be made later and not
here: clip (overflow saturates at $1$), project (renormalize
direction and forget extra strength), or refuse (the operator is
undefined off the ball). Schoolbook $+$ on the $\mathtt{num}$
encoding is a different door: decode, add in $\mathbb{R}$, encode
back if the chart allows. Do not let $\mathbb{R}^3$ addition
quietly replace $\oplus$, $\boxplus$, or $T$. The embedding
lives in $\mathbb{R}^3$. The unit lives in the ball.
As already noted in Section~\ref{subsec:register_scope_final},
the three-dimensional ball teaches the algebra; close document
semantics, if pursued, live on a higher-dimensional sphere that
carries the same verbs.

\paragraph{Dimensionality.}
\label{subsec:dimensionality_open}
Nothing stated
above is intrinsically three-dimensional; the stage may be
$S^{N_{\mathrm{dim}}-1}$ as well as $S^2$. The magnitude cap
$r\in[0,1]$ is a property of a vector's length, not of the ambient
dimension. The log and exponential maps used for the offset operator
(Equations around~\eqref{eq:sphere_exp}--\eqref{eq:analogical_offset})
are already Riemannian statements that hold on any $S^{N_{\mathrm{dim}}-1}$ in
$\mathbb{R}^{N_{\mathrm{dim}}}$, not only $S^2$. $\mathrm{SO}(3)$ generalizes to
$\mathrm{SO}(N_{\mathrm{dim}})$ as an abstract group --- associativity and
invertibility are group-theoretic facts about $\mathrm{SO}(N_{\mathrm{dim}})$ for
any $N_{\mathrm{dim}}$, not special to $N_{\mathrm{dim}}=3$ --- but that is a statement about the
group, not about $T_w$'s specific construction from one vector:
\S\ref{subsec:transform_higher_dim} found that a single $\hat n_w$
does not carry enough data to fix a rotation outside three
dimensions, however true the group's own properties remain. Three
dimensions were chosen for legibility and hand-checkable worked
examples, not necessity. Raising $N_{\mathrm{dim}}$ is expected to relieve one
specific problem: too few directions forcing distinct concepts to
crowd onto nearby points, a kind of clipping by dimensional scarcity
rather than by magnitude. Because crowding is a many-atom problem, the
benefit is conditional on the bowl actually being large enough to
crowd; raising $N_{\mathrm{dim}}$ for a bowl of two or three atoms relieves nothing,
since ordinary distinct points already sit far apart on $S^2$. It does
\emph{not} relieve the separate
overflow problem discussed just above (clip, project, or refuse when
$r$ exceeds 1) --- that is a magnitude question and stays open
regardless of $N_{\mathrm{dim}}$.

$T$ and offset are no longer part of this TODO.
\S\ref{subsec:transform_higher_dim} restates $T$ for general
$N_{\mathrm{dim}}$, invertible and non-commutative up to
$N_{\mathrm{dim}}=50$ and exact at $N_{\mathrm{dim}}=3$;
\S\ref{subsec:offset_logexp} restates offset's transport as a
rotation within the plane of its two seats, exact at
$N_{\mathrm{dim}}=3$. The remaining operators contain no step that
depends on the dimension. Gather is set union. Mix reads phase among
units that already share a direction. Assembly's bracket
$3+\hat n_A\cdot\hat n_B$ stays within $[2,4]$ in any dimension, so it
never repels by the same argument as in three. The identities of
\hyperref[subsec:axiom_v1_metrics]{Axiom~7} are dot-product
identities. Each follows from the argument just given, in any
dimension; no run is needed to establish it. What stays three-dimensional is presentation: the
polar chart $(\theta,\phi)$ and every worked example, by the
convention stated at the start of this paper. Mix across differing
directions remains open (\S\ref{subsec:mix_scattered_open}), but
that is a question about direction, not dimension.

The algebra is dimension-blind and count-honest: every informational
arrow is a unit; the stage may be $S^2$ or $S^{N_{\mathrm{dim}}-1}$;
packing more arrows into one 3-D atom is a density choice, not a
change of claim. 3-D with a high $N$ is the right teaching object.
Higher $N_{\mathrm{dim}}$ is the same object with more sky --- use it
when crowding, not when the idea feels small.

\paragraph{Open: Mix across differing directions.}
\label{subsec:mix_scattered_open}
\hyperref[subsec:axiom_interference]{Axiom~5} defines Mix only among states
that share a direction $\hat n$, where it is phasor combination on
$\chi$. What Mix should mean for states at different directions is not
settled here. An earlier draft of this paper carried a provisional
recipe --- normalize the vector sum $r_A\hat n_A+r_B\hat n_B$ --- and
it is withdrawn, because it is the only place in the construction
where $\hat n$ would move by something other than a rotation.
Everywhere else, direction is the exclusive property of
$\mathrm{SO}(3)$ (\S\ref{subsec:axiom_rotation}) and of $T$, and a
blend that slides an arrow between axes is not a rotation. Three
routes remain open, and this paper picks none of them: refuse Mix
across directions outright, making the operator partial in place just
as it is already partial in sort; rotate one state onto the other's
axis by an explicit, stated $T$ first, so that all blending stays
phasor blending; or accept a genuine two-stage operator and state
separately what it preserves. Whichever is chosen, it should be shown
that the resulting map still respects the cap and does not smuggle
$\mathrm{SO}(3)$'s job into $\boxplus$.

\paragraph{Open: reducing a path by merging close actors.} Two
\emph{sequential} actors already compose into one equivalent rotation
for free --- that is just $\mathrm{SO}(3)$'s (or $\mathrm{SO}(N_{\mathrm{dim}})$'s)
closure, already latent in $T$'s associativity. That is not the open
question. The open question is whether several actors that are
merely \emph{close} --- by some stated distance on $\mathcal{S}$, not
identical --- may be replaced by one representative actor before a
path is walked, trading path length for approximation. This is a
coarse-graining or clustering step, not composition, and it is not
defined anywhere in this paper. Two things would need stating before
it could be trusted: a distance and a merge rule (is the
representative a $\boxplus$-blend of the close actors, an average
rotation, or something else), and a bound relating that distance to
how far the reduced path's endpoint may drift from the unreduced
path's endpoint. Composition is associative and exact;
this proposed reduction is neither, and its error, if any, has not
been characterized.
A merge of this kind, whatever form it eventually takes, would
produce an additional, derived actor for a shorter path to use --- it
would not remove the original close actors from the bowl. The bowl
under $\oplus$ (Section~\ref{subsec:process_principle}) already keeps
every distinct seated unit; a reduced actor sits alongside its
sources, not in place of them, and the unreduced path through the
original actors remains walkable and exact. Reduction would be a
second, optional reading of the same bowl, the way Mix and $T$
already are, not a rewrite of what the bowl holds. Left open.

\paragraph{Open: location as a sort.} A position (latitude and
longitude, say, or any other two-angle coordinate) is not a fifth
coordinate on every atom; the atom stays the 4-tuple of
\hyperref[subsec:axiom1]{Axiom~1}, and $\hat n$ remains a pattern's semantic
address, not its physical one. A location instead sits for admission
the way $\mathtt{chladni}$ did
(Section~\ref{subsec:chladni_inclusion}): its own $\tau$, its own
chart, and its own $\varphi_\tau$ stating what native law on positions
$\boxplus$ or $T$ would have to honor. Neither has been attempted
here. Two candidates worth weighing when it is: whether combining two
locations behaves enough like $\oplus$'s union of presence to admit at
Tier 0 only --- plausible if the native operation absorbs (a smaller
region inside a larger one yielding the larger), since $\oplus$ has no
absorption law and cannot honor one; or whether some native law on
positions maps onto $T$ or $\boxplus$ at Tier 1, the way
$\mathtt{chladni}$'s modal superposition did. Left open.

This is a separate question from the seat $\vec P$ of
Section~\ref{subsec:enclosure_seats}, and admitting one would not
settle the other. A seat is external bookkeeping --- a place a unit
occupies, carrying no meaning and entering no chart. A location
\emph{sort} would be a unit in its own right, seated in bowls
alongside temperature and wind, and would have to pass Axiom~0 on its
own terms before $\boxplus$ or $T$ could touch it.

\end{document}